\documentclass[11pt,a4paper]{article}
\usepackage[T1]{fontenc}
\usepackage{lmodern}
\ifdefined\pdfgentounicode
  \input{glyphtounicode}
  \pdfglyphtounicode{parenleftbig}{0028}
  \pdfglyphtounicode{parenrightbig}{0029}
  \pdfglyphtounicode{parenleftbigg}{0028}
  \pdfglyphtounicode{parenrightbigg}{0029}
  \pdfglyphtounicode{bracketleftbigg}{005B}
  \pdfglyphtounicode{bracketrightbigg}{005D}
  \pdfglyphtounicode{vextendsingle}{007C}
  \pdfglyphtounicode{vextenddouble}{2016}
  \pdfgentounicode=1
\fi
\usepackage[margin=27mm,headheight=15pt]{geometry}
\usepackage{amsmath,amssymb,amsthm,mathtools}
\usepackage{booktabs,longtable,array,enumitem}
\usepackage{microtype,xcolor,fancyhdr}
\usepackage[colorlinks=true,linkcolor=blue!45!black,citecolor=blue!45!black,urlcolor=blue!45!black,breaklinks=true]{hyperref}
\hypersetup{pdftitle={A compensated-process approach to the no-three-in-line problem},pdfauthor={},pdfsubject={English research proof draft; conditional numerical target 1.501n}}
\setlength{\emergencystretch}{2em}
\setlength{\parskip}{3pt}
\setlist{itemsep=2pt,topsep=4pt}
\newtheorem{theorem}{Theorem}[section]
\newtheorem{lemma}[theorem]{Lemma}
\newtheorem{proposition}[theorem]{Proposition}
\newtheorem{corollary}[theorem]{Corollary}
\theoremstyle{definition}
\newtheorem{definition}[theorem]{Definition}
\newtheorem{remark}[theorem]{Remark}
\newcommand{\E}{\mathbb E}
\newcommand{\Prob}{\mathbb P}
\newcommand{\F}{\mathcal F}
\newcommand{\R}{\mathbb R}
\newcommand{\N}{\mathbb N}
\newcommand{\Z}{\mathbb Z}
\newcommand{\ind}{\mathbf 1}
\newcommand{\dd}{\,\mathrm d}
\newcommand{\Var}{\operatorname{Var}}
\newcommand{\NoThree}{\operatorname{NoThree}}
\newcommand{\supp}{\operatorname{supp}}
\newcommand{\col}{\operatorname{col}}
\newcommand{\typ}{\operatorname{typ}}
\newcommand{\pos}[1]{\left[#1\right]_+}
\newcommand{\given}{\,\middle|\,}
\newcommand{\repo}{https://github.com/agolkjaoi/no-three-in-line-lean}
\pagestyle{fancy}
\fancyhf{}
\fancyhead[L]{\small No-three-in-line: a compensated-process proof draft}
\fancyhead[R]{\small 10 November 2026}
\fancyfoot[C]{\thepage}
\renewcommand{\headrulewidth}{0.3pt}
\title{\textbf{A compensated-process approach\\to the no-three-in-line problem}\\[5pt]
\large An English proof draft for the target \(1.501n\)}
\author{}
\date{10 November 2026\\[3pt]\small Version 0.1: mathematical exposition with explicit verification obligations}
\begin{document}
\maketitle
\begin{abstract}
We describe a finite, continuous-time construction of sets containing no three
collinear points in the \(n\times n\) integer grid. Birth rates are compensated
for geometric deletion and for the two available places on selected families
of lines. A saved-credit martingale gives an exact identity: the expected
output equals deterministic input mass minus the remaining exposure lost at
early retirement. This identity separates a small fixed error budget from
the geometry and concentration estimates needed to meet it.

The manuscript develops the lattice estimates, mixed-moment argument,
protected-line concentration, long-line and short-star decomposition,
finite sieve, and rational geometric checking behind this reduction.
It incorporates the two followups to the original preparation notes.
The numerical application is explicitly conditional: complete checking and
matching of the candidate input, including its mass and root-deadline pair
bounds, are still required. The accompanying partial Lean development does
not prove the unconditional \(1.501n\) theorem. This is a proof draft for
mathematical examination, not a claim that the target has been established.
\end{abstract}
\noindent\textbf{Reading and attribution.}
No knowledge of Lean is needed for the mathematical text.
The detailed stochastic arguments are research arguments under review;
their inclusion here is not a new independent verification.
AI systems contributed substantially to the research, code, proofs, checking,
and this exposition. No named author is supplied. Precise disclosure and
verification scope appear in Section~\ref{sec:status}.
\tableofcontents
\newpage
\section{The target and the structure of the argument}\label{sec:overview}

Write \(V_n=\{0,\ldots,n-1\}^2\). A set \(S\subseteq V_n\) has the
\emph{no-three-in-line property} if every three distinct points of \(S\)
are noncollinear. Equivalently,
\[
 (q_x-p_x)(r_y-p_y)-(q_y-p_y)(r_x-p_x)\ne0
 \qquad(p,q,r\in S\text{ distinct}).
\]
Every such set has at most two points on each horizontal row, hence
\(|S|\le2n\). The intended numerical conclusion is
\begin{equation}\label{eq:target}
 \exists n_0\ \forall n\ge n_0\ \exists S\subseteq V_n:
 \quad \NoThree(S),\qquad 1501n\le1000|S|.
\end{equation}
The positive improvement over \(3n/2\) is the fixed number \(n/1000\).
The quantifier is over every sufficiently large integer \(n\), with no
divisibility restriction.

\subsection{Finite input, fixed before the grid size}
Choose a finite rational time partition
\(0=t_0<t_1<\cdots<t_J=T\) and a fixed finite rational rectangular
partition of the unit square. A common refinement may be taken to be a
\(G\times G\) grid for a fixed integer \(G\). Let \(\mathcal B\) be the
union of its horizontal and vertical grid lines.
There are four parity types \(r\in(\Z/2\Z)^2\). For each type choose
a deterministic function
\[
 f_{r,t}:[0,1]^2\longrightarrow[0,1]
\]
that is constant on every open spatial cell and every time interval.
Set it to zero on \(\mathcal B\), outside the square, and after time \(T\).
Define
\begin{equation}\label{eq:input}
 m_t(p)=f_{\typ(p),t}(p/n),\qquad
 F_{r,t}(x)=\int_0^t f_{r,u}(x)\dd u,\qquad
 A=\frac14\sum_r\int_{[0,1]^2}F_{r,T}(x)\dd x.
\end{equation}
Here \(\typ(p)=p\bmod2\). The factor \(1/4\) belongs to the spatial
mass, not to the birth rate of an individual point.
The last-use time \(\tau_r(x)\) is the right endpoint of the last interval
on which that type and cell is used; it is zero for an unused cell.
For a grid point write \(\tau_p=\tau_{\typ(p)}(p/n)\).

In the candidate data the spatial mask in a phase is obtained from a
union \(U\) of accepted closed rectangles and their prescribed images:
the actual region is \(\operatorname{int}(U)\setminus\mathcal B\).
Areas are union areas. Overlapping images are never added with
multiplicity. Removing \(\mathcal B\) is part of the definition of the
discrete process, even though it has zero planar area.

\subsection{The geometric input conditions}
A primitive integer direction \(v\) is identified with \(-v\). Put
\(h(v)=\max(|v_1|,|v_2|)\) and \(\bar v=v\bmod2\ne0\).
For a field \(F\) define the major-coordinate chord integral
\[
 C_vF(x)=\int_{\R}F\bigl(x+s v/h(v)\bigr)\dd s.
\]
This is not an arc-length integral. Let \(\Gamma\) be a fixed finite
family of designated directions containing both coordinate directions,
of cardinality \(g\).
The numerical example uses all directions of height at most three,
so \(g=16\). Define the line capacity
\begin{equation}\label{eq:capacity}
 L_{v,r}(F;x)=\frac{C_v(F_r+F_{r+\bar v})(x)}{2h(v)}.
\end{equation}
Let the four direction families be
\((1,u),(1,-u),(u,1),(-u,1)\), \(0\le u\le1\), and let
\(c_{F,j}(x,u)\) denote the corresponding chord. Set
\begin{equation}\label{eq:risk}
 g_F(x)=\frac14\sum_{j=1}^4\int_0^1c_{F,j}(x,u)^2\dd u,
 \qquad
 \mathcal R_r(F;x)=\frac{D_*}{6}\sum_{s\ne r}g_{F_r+F_s}(x).
\end{equation}
Section~\ref{sec:sieve} gives a finite arithmetic criterion for a valid
coefficient \(D_*\); the example uses \(76/125\).

The finite schedule is \emph{admissible} for fixed parameters if there
are positive constants \(\kappa_H,\kappa_L\) and a number \(0\le r_*<d\)
such that
\begin{align}
 L_{v,r}(F_T;x)&\le K-\kappa_H &&(v\in\Gamma),\label{eq:heavyinput}\\
 L_{v,r}(F_T;x)&\le\lambda-\kappa_L &&(v\notin\Gamma),\label{eq:lightinput}\\
 \mathcal R_r(F_{\tau_r(x)};x)&\le r_* &&(\tau_r(x)>0),\label{eq:riskinput}
\end{align}
uniformly over the relevant roots, types and lines. In the last formula
\emph{every one of the four fields is evaluated at the same root's
deadline}. The constants satisfy
\begin{equation}\label{eq:exponents}
 \begin{gathered}
 0<K<2,\quad 0<\rho<1,\quad 0\le\lambda<1-\rho,\\
 0<d<\gamma<1,\quad 0<\xi<d,\quad \zeta>0,\quad 1-d-2\xi-\zeta>0.
 \end{gathered}
\end{equation}
All data and all displayed constants are fixed before \(n\) tends to
infinity. Constants depending on a moment order may depend on that
order, but the order is also fixed before taking the limit.

\begin{theorem}[Analytic reduction, in the scope of this proof draft]\label{thm:bridge}
For an admissible finite schedule, the stopped compensated process
defined in Section~\ref{sec:process} produces a random set \(S_n\subseteq
V_n\) with the no-three-in-line property and
\[
 \E|S_n|=(A-o(1))n.
\]
Consequently, for every fixed \(c<A\), a set of size at least \(cn\)
exists for every sufficiently large \(n\).
\end{theorem}

The detailed argument for this theorem occupies
Sections~\ref{sec:geometry}--\ref{sec:closing}. The numerical use of the
theorem is a separate obligation: the particular finite schedule must
actually satisfy \eqref{eq:heavyinput}--\eqref{eq:riskinput}, with the
claimed union mass and chronological matching. This manuscript does
not certify that entire concrete input. In particular, the theorem
is not an unconditional assertion of \eqref{eq:target}.

\subsection{Why saved credit is the useful quantity}
Each active point has a weight \(w_p\), initially one. It is selected at
rate \(m_t(p)w_p/n\). Weights grow to compensate for the chance of
selection, deletion by a newly completed pair, or a prescribed
halving. At an early safety stop, the positive weight is saved rather
than discarded. The resulting process \(Z_p\) is a bounded martingale
with \(\E Z_p(t)=1\).

The change of probability measure
\[
 \dd\Prob^{(p)}=Z_p(T)\dd\Prob
\]
measures the amount of credit lost by retiring \(p\), rather than just
the raw probability that \(p\) retires. Let \(\theta_p\) be its first
positive saving time and \(E_p\) the event that this happens strictly
before \(\tau_p\) for a safety reason. Put
\[
 A_n=\frac1{n^2}\sum_{p\in V_n}F_{\typ(p),T}(p/n).
\]
The exact identity proved in Section~\ref{sec:loss} is
\begin{equation}\label{eq:previewloss}
 A_n-\frac{\E|S_n|}{n}
 =\frac1{n^2}\sum_p\E^{(p)}\left[
 \ind_{E_p}\int_{\theta_p}^{T}m_t(p)\dd t\right].
\end{equation}
Thus a bound \(\sup_p\Prob^{(p)}(E_p)\le\varepsilon\) implies
\(\E|S_n|/n\ge(1-\varepsilon)A_n\).
There are five causes to control: weight cap, ordinary-line tag,
pair-correction floor, virtual-deletion bound, and designated-line clock.
The last three are global safety stops. The paper proves their estimates
on the already stopped process; it does not assume that those stops
are rare while deriving the estimates that make them rare.

\subsection{A fixed budget suffices for the numerical target}
The recorded scaled mass and target are
\begin{equation}\label{eq:massnumbers}
 A_* =\frac{5036868114751900569}{3355443200000000000},\qquad
 c=\frac{1501}{1000},\qquad
 \Delta=A_*-c=\frac{347871551900569}{3355443200000000000}.
\end{equation}
These are candidate input values; verifying the mass of the accepted
union remains part of the concrete application. If they are justified,
and if \(A_n\ge A_*-\Delta/4\) and
\(\sup_p\Prob^{(p)}(E_p)\le10^{-5}\), then
\[
 (1-10^{-5})(A_*-\Delta/4)-c
 =0.0000627445734908543514408171176910400390625>0.
\]
This strict rational inequality is enough for \eqref{eq:target}.
There is no need to prove an effective, practically small threshold
\(n_0\), or to compute an unbounded sequence of ever tighter error
budgets. The asymptotic argument is a way to obtain one fixed budget.


\section{Uniform lattice geometry}\label{sec:geometry}

In this section a line means its intersection with \(V_n\), and only
lines containing at least two grid points are counted. Write \(L_\ell\)
for its cardinality and \(L_{pq}\) for the length of the line through
distinct \(p,q\). Constants in this section are uniform in the root's
position, including roots approaching a cell boundary as \(n\) changes.

\begin{lemma}[Direction and line counts]\label{lem:lines}
There are at most \(4h\) primitive unoriented directions of height
\(h\). Every grid line of that height has at most \(2n/h\) points.
Consequently, uniformly in \(p\),
\begin{align}
 \sum_{\ell\ni p}L_\ell^2&\le16n^2(1+\log n),\label{eq:line2}\\
 \sum_{\ell\ni p}L_\ell^k&\le
 4(2n)^k\left(1+\frac1{k-2}\right)\qquad(k>2).
 \label{eq:linek}
\end{align}
\end{lemma}
\begin{proof}
If \(v\) is primitive, the integer points on its line are precisely
\(p+\Z v\). Indeed, Bezout coefficients show that a real multiple of
\(v\) with integer coordinates has an integer multiplier. A coordinate
of absolute step \(h\) gives \(L_\ell\le1+(n-1)/h\le2n/h\).
For \(h\ge2\), counting the primitive points on the boundary of the
square of radius \(h\), with opposite signs identified, gives
\(4\varphi(h)\le4h\); height one has four directions. Sum
\(4h(2n/h)^k\), using the harmonic bound when \(k=2\) and an integral
comparison when \(k>2\). The same argument gives
\(\sum_{q\ne p}L_{pq}=\sum_{\ell\ni p}(L_\ell-1)L_\ell
\le\sum L_\ell^2\).
\end{proof}

\begin{lemma}[Two uniform incidence kernels]\label{lem:kernels}
For every fixed \(\xi>0\),
\begin{align}
 \sum_{q\notin\ell(a,b)}L_{aq}L_{bq}&\le C_\xi n^{2+\xi}
       &&(a\ne b),\label{eq:twostar}\\
 \sum_{q\in\ell\cap V_n}L_{pq}&\le C_\xi n^{1+\xi}
       &&(p\notin\ell).
       \label{eq:fixedline}
\end{align}
\end{lemma}
\begin{proof}
The elementary divisor estimate \(\tau(m)\le C_\epsilon m^\epsilon\)
follows by writing \(m=\prod l^{a_l}\). For sufficiently large primes,
\(a+1\le2^a\le l^{\epsilon a}\); the finitely many smaller primes have
finite suprema of \((a+1)/l^{\epsilon a}\). Hence, with
\(\Delta_n=\max_{m\le2n^2}\tau(m)\),
\(\Delta_n(1+\log n)^2\le C_\xi n^\xi\).

For \eqref{eq:twostar}, write \(q-a=e u\), \(q-b=f v\), where
\(e,f>0\) and \(u,v\) are primitive, and put \(w=b-a\).
Then \(eu-fv=w\) and \(f\mid\det(u,w)\). The exclusion of the baseline
ensures \(1\le|\det(u,w)|\le2n^2\), so for fixed \(u\) there are at
most \(\Delta_n\) choices of \(f\). In dyadic height bands
\(H\le h(u)<2H\), \(J\le h(v)<2J\), assume \(H\le J\).
There are \(O(H^2)\) choices of \(u\). The congruence
\(eu\equiv w\pmod f\) confines \(e\) to at most one residue class
modulo \(f\), by Bezout. A coordinate with \(|u_i|\ge H\), together
with \(|v_i|<2J\), confines \(e/f\) to an interval of length at most
\(4J/H\). Thus the number of candidates is
\(O(\Delta_n HJ)\). Their individual weights are at most
\(4n^2/(HJ)\). Interchange \(a,b\) when \(H>J\), then sum the
\(O((1+\log n)^2)\) bands.

For \eqref{eq:fixedline}, the cases of zero or one grid point are
immediate. Otherwise write \(\ell\cap\Z^2=a+\Z v\) with primitive
\(v\), and let \(D=|\det(v,p-a)|\in[1,2n^2]\).
For \(q\in\ell\), put \(g(q)=\gcd(|q_1-p_1|,|q_2-p_2|)\).
Then \(g(q)\mid D\). For each \(e\mid D\), the condition
\(q\equiv p\pmod e\) restricts the line parameter to one residue class.
Within \(\|q-p\|_\infty\le R\), this gives at most
\(2R/(e h(v))+1\) choices. Since
\(g(q)\le\sum_{e\mid g(q)}e\) and relevant \(e\le R\),
\[
 \sum_{\|q-p\|_\infty\le R}g(q)\le3R\tau(D).
\]
In the shell \(R/2<\|q-p\|_\infty\le R\),
\(L_{pq}\le4ng(q)/R\), so the shell contributes at most
\(12n\tau(D)\). Sum dyadic shells and use the divisor estimate.
\end{proof}

\subsection{Parity, chords and a uniform boundary estimate}
Suppose \(F_r\) are constant on a fixed partition into \(R\) open
rectangular cells, vanish on the common boundaries, and lie in
\([0,M]\). Along a primitive line through a root of type \(r\), the
types alternate between \(r\) and \(r+\bar v\).

\begin{lemma}[Lattice-to-chord comparison]\label{lem:chord}
For every root and primitive direction of height \(h\),
\begin{equation}\label{eq:chorderror}
 \left|\sum_{q\in\ell(p,v)\setminus\{p\}}
 F_{\typ(q)}(q/n)
 -\frac n{2h}C_v(F_r+F_{r+\bar v})(p/n)\right|
 \le(4R+1)M.
\end{equation}
Moreover \(A_n=A+O_{R,T}(n^{-1})\).
\end{lemma}
\begin{proof}
Each parity subsequence has major-coordinate spacing \(2h/n\).
The number of points of any shifted one-dimensional lattice in an
open interval differs from interval length divided by spacing by at
most two. Sum over the two types and the \(R\) cells, then remove the
root, costing at most \(M\). A line coinciding with a common boundary
has field identically zero on both sides of \eqref{eq:chorderror}.
For the mass statement, count each coordinate of each parity class
in an open rectangle: the count is
\(n^2\operatorname{area}(E)/4+O(n)\). Sum the fixed number of cells
and types. No assumption that \(n\) is even or a multiple of \(G\)
is used.
\end{proof}

\begin{lemma}[Uniform bounded variation]\label{lem:bv}
For a nonnegative field \(F\le M\) on \(R\) fixed rectangular cells,
each of the four functions \(u\mapsto c_{F,j}(x,u)^2\) satisfies
\[
 \|c_{F,j}(x,\cdot)^2\|_\infty\le M^2,
 \qquad \Var(c_{F,j}(x,\cdot)^2)\le(4R+2)M^2,
\]
uniformly in \(x\in[0,1]^2\).
\end{lemma}
\begin{proof}
For \((1,u)\), the chord of a rectangle \((a,b)\times(c,d)\) through
\((x,y)\), for \(u>0\), is
\[
 \pos{\min(b,x+(d-y)/u)-\max(a,x+(c-y)/u)}.
\]
Its changes of branch occur only at slopes from the root to the four
vertices. Commonly refine the at most \(4R\) switch values. On each
interval the total chord has the form \(\alpha+\beta/u\), hence is
monotone, as is its square since it is nonnegative. The square lies
in \([0,M^2]\). At positive switches the values agree. At zero, a line
may coincide with a common boundary; its possible endpoint jump costs
at most one further \(M^2\). The stated bound follows. Reflections
and coordinate exchange give the other three families.
\end{proof}

This proof uses bounded variation, not a uniform derivative bound.
For example, for a half-square and a root a distance \(\epsilon\)
above its boundary, a chord can have the form
\(\max(0,1/2-\epsilon/u)\). Its derivative near \(u=2\epsilon\)
is unbounded as \(\epsilon\downarrow0\), whereas its squared
variation remains at most \(1/4\).

\subsection{A finite sieve for primitive directions}\label{sec:sieve}
Let \(P\) be a fixed finite set of pairwise coprime odd integers greater
than one; primality is unnecessary. Put
\[
 Q=\prod_{p\in P}p,\qquad
 D_P=\frac34\prod_{p\in P}(1-p^{-2}),\qquad
 A_P=\prod_{p\in P}(1+p^{-2}).
\]
For \(e=(\epsilon,\eta)\in\{(0,1),(1,0),(1,1)\}\) and a nonnegative
BV function \(w\) on \([0,1]\), define
\[
 K_{N,e}(w)=\sum_{h\le N}\frac1{h^2}
 \sum_{\substack{1\le a\le h\,;\ (a,h)=1\\(h,a)\bmod2=e}}w(a/h).
\]
Write \(I=\int_0^1w\), \(M_w=\sup w\), and \(V_w=\Var(w)\).

\begin{lemma}[Finite-sieve upper bound]\label{lem:sieve}
For every real \(N\ge Q\),
\begin{equation}\label{eq:sieve}
 K_{N,e}(w)\le\frac{D_P}{3}I\log N+C_P(M_w+V_w),\qquad
 C_P=\left(5+\frac14\log Q\right)A_P.
\end{equation}
The rational upper bound
\(C_P\le[5+\tfrac14\sum_{p\in P}(p-1)]A_P\) is also valid.
\end{lemma}
\begin{proof}
For real \(x\ge1\), omit the coprimality condition and denote the
corresponding sum by \(J_e(x;w)\). Counting alternating tags gives
\[
 \left|\sum_{\substack{1\le b\le k\\b\equiv\eta\ (2)}}w(b/k)
       -\frac k2 I\right|\le2M_w+V_w.
\]
To see this, place intervals of length \(2/k\) to the left of the
tags; the oscillation errors sum to at most \(V_w\), and the omitted
or cut end intervals have total length at most \(4/k\).
Also the parity harmonic sum differs from \(\tfrac12\log x\)
by at most two, by comparing ordinary harmonic sums with integrals.
The finite estimate \(\sum_{k\le m}k^{-2}\le2\) therefore gives
\[
 J_e(x;w)=\frac I4\log x+E_e(x),\qquad
 |E_e(x)|\le5(M_w+V_w).
\]
Nonnegativity of \(w\) allows the bound
\[
 \ind_{(a,h)=1}\le
 \prod_{p\in P}(1-\ind_{p\mid a,\ p\mid h}).
\]
Expand this finite product. For \(A\subseteq P\), put
\(d_A=\prod_{p\in A}p\) and \(s_A=(-1)^{|A|}\).
Pairwise coprimality identifies simultaneous divisibility with
divisibility by \(d_A\); oddness preserves the parity class. Thus
\[
 K_{N,e}(w)\le\sum_{A\subseteq P}\frac{s_A}{d_A^2}J_e(N/d_A;w).
\]
Since \(N\ge Q\), every \(N/d_A\ge1\). Substituting the two-sided
estimate for \(J_e\) and using
\(\sum s_A/d_A^2=4D_P/3\), \(\sum1/d_A^2=A_P\),
\(\log d_A\le\log Q\), and \(I\le M_w\) proves
\eqref{eq:sieve}. It is essential to use a two-sided error estimate
inside the signed inclusion-exclusion sum. Finally \(\log p\le p-1\)
gives the rational constant.
\end{proof}

Only a one-sided bound is needed in this application. Consequently
no infinite Mobius series, value of \(\zeta(2)\), or assertion of
probabilistic independence between parity and direction is required.
The recorded choice of 175 factors gives \(D_P<76/125\); its product
\(Q\) has 440 decimal digits. This large fixed threshold is compatible
with the asymptotic quantifiers, but gives no useful small-grid claim.
The factors and rational comparison are included in the companion data.

\subsection{The common-deadline pair sum}
For four fields \(F\) and a root \(p\), define
\begin{equation}\label{eq:pairsum}
 \mathcal B_n(p;F)=\frac1{2n^2}
 \sum_{\substack{q,z\in V_n\setminus\{p\},\ q\ne z\\p,q,z\text{ collinear}}}
 F_{\typ(q)}(q/n)F_{\typ(z)}(z/n).
\end{equation}
The pair sum is ordered; the factor \(1/2\) removes this duplication.
A superscript \(\Gamma^c\) restricts it to nondesignated directions.

\begin{proposition}[Uniform risk comparison]\label{prop:pair}
If \(D_*\ge D_P\), then for \(n\ge\max(2,Q)\),
\[
 \mathcal B_n^{\Gamma^c}(p;F)
 \le\mathcal R_{\typ(p)}(F;p/n)\log n+C_{R,M,P},
\]
where the constant is independent of \(n,p\) and of the values of the
fields within their fixed partition and bound.
\end{proposition}
\begin{proof}
For each direction let \(W_v\) be the weighted line sum with the root
removed. The contribution is
\((W_v^2-\sum_qF_{\typ(q)}(q/n)^2)/(2n^2)\).
By \eqref{eq:chorderror},
\(W_v=nC_v(F_r+F_{r+\bar v})/(2h)+e_v\),
\(|e_v|\le E_0=(4R+1)M\). Expanding the square and using
at most \(4h\) directions at height \(h\), all cross errors and
diagonal terms sum to \(O_{R,M}(1)\) after normalization. Thus
\[
 \mathcal B_n(p;F)=\frac18\sum_{\substack{v\text{ primitive}/\pm\\h(v)\le n}}
 \frac{C_v(F_r+F_{r+\bar v})(p/n)^2}{h(v)^2}+O_{R,M}(1).
\]
Represent nonaxis directions by the four families with slope
\(a/h\), \(1\le a\le h\). The diagonal overlap can be retained for
an upper bound; the two axes contribute a bounded amount separately.
Apply Lemma~\ref{lem:sieve} to each of the twelve family/parity
functions, whose BV norms are uniform by Lemma~\ref{lem:bv}.
For each nonzero parity \(v\), the four integrals sum to
\(4g_{F_r+F_{r+v}}\). The coefficient is
\(\tfrac18\cdot\tfrac{D_P}{3}\cdot4=D_P/6\).
Increase it to \(D_*/6\) and discard the nonnegative designated
contributions.
\end{proof}

For the time-dependent fields take \(M=T\) once and for all. The same
constant works for every time, including every root's last-use time.
For a fixed root with deadline \(\tau\), finite-product differentiation
gives the exact identity
\begin{equation}\label{eq:timesym}
 \frac1{n^2}\int_0^\tau\sum_{q,z}^{\Gamma^c}m_t(q)F_{\typ(z),t}(z/n)\dd t
 =\mathcal B_n^{\Gamma^c}(p;F_\tau).
\end{equation}
Indeed, differentiate each product \(F_t(q)F_t(z)\), exchange \(q,z\),
and integrate. The index set is symmetric and consists of distinct
points different from the root. Using separate deadlines for \(q\)
and \(z\) would not give this identity.


\section{The stopped compensated process}\label{sec:process}

Fix \(n\ge2\) and put \(B=n^d\), \(M_n=n^\gamma\), and
\(\epsilon_n=n^{-\zeta}\). A point is active, selected, geometrically
dead, or positively retired. Initially every point is active, its weight
is one, and \(S=\varnothing\). Selection is permanent. Dead and retired
points never return to the active set \(Q_t\).

\subsection{Line variables, rates and flows}
Every designated line has occupancy \(N_\ell=|S\cap\ell|\) and clock
\(k_\ell\). For \(N=0,1\) set
\[
 a_N(k)=\frac{2-N}{2-k},\qquad b_0=\frac12,\quad b_1=1.
\]
Every nondesignated line has a credit \(r_\ell\), initially one.
It changes only while the line contains exactly one selected point.
Before the global stopping time \(\sigma\), define
\begin{align}
 \lambda_p(t)&=\ind_{p\in Q_t}\frac{m_t(p)w_p(t)}n,
 &\Lambda_\ell&=\sum_{p\in\ell}\lambda_p,\label{eq:rates}\\
 \dot k_\ell&=\Lambda_\ell/a_{N_\ell}(k_\ell)
       &&(N_\ell=0,1),\label{eq:clockflow}\\
 \dot r_\ell&=-r_\ell\Lambda_\ell
       &&(N_\ell=1,\ \ell\text{ nondesignated}).\label{eq:tagflow}
\end{align}
The other clock and line-credit derivatives are zero. For active \(p\),
put
\[
 H_p=\sum_{q\in Q_t\setminus\{p\}}\lambda_q
 \ind_{\ell(p,q)\notin\Gamma,\ |S\cap\ell(p,q)|=1}.
\]
Here \(\ell\notin\Gamma\) means that its direction is nondesignated.
The weight flow is
\begin{equation}\label{eq:weightflow}
 \frac{\dot w_p}{w_p}=H_p+\lambda_p+
 \sum_{\substack{\ell\ni p\\\ell\in\Gamma}}
 (\Lambda_\ell-\lambda_p)b_{N_\ell}.
\end{equation}
The term \(H_p\) counts each actual cause \(q\) once, not once for
each possible witness. All rates are set to zero after \(\sigma\).

At a birth of \(q\), use the pre-birth state to choose the branch,
then perform these operations in order:
\begin{enumerate}
\item Select \(q\) and set its weight to zero.
\item Kill every active point that now lies on a line with two selected
points, setting its weight to zero.
\item On each designated line having its first birth, halve the weights
of the remaining active points.
\item Update occupancies and the activation flags of ordinary lines.
The values of \(k_\ell\) and \(r_\ell\) are not reset.
\item Process any newly reached safety boundary and positive retirement.
\end{enumerate}
For a fixed other point \(p\), at most one line through \(p,q\) can
cause a halving. Geometric death takes priority over positive saving.

\subsection{Safety boundaries and the observation convention}
A point retires positively at its first weight-cap hit \(w_p=B\),
at a hit \(r_\ell=\rho\) on an ordinary line through it, or at its
deadline \(\tau_p\). Its current positive weight is saved. Unused
points retire at time zero with saved weight one.

Two additional variables define global guards. For distinct active
\(a,b\), let \(\kappa_{ab}\) be the sum of rates of birth causes
\(q\notin\ell(a,b)\) that would change both weights by a factor less
than one. Set \(\dot\eta_{ab}=-\kappa_{ab}\eta_{ab}\), initially
\(\eta_{ab}=1\), and freeze this variable when either endpoint leaves
the active state. For a designated line \(\ell\) and an active point
\(q\notin\ell\), define the virtual deletion mass
\begin{equation}\label{eq:virtual}
 \mathcal V_{\ell,q}=
 \sum_{y\in\ell\cap Q_t}\frac{w_y}{a_{N_\ell}(k_\ell)}
 \ind_{\ell(q,y)\notin\Gamma,\ |S\cap\ell(q,y)|=1}.
\end{equation}
Globally stop at the first hit of \(k_\ell=K\),
\(\eta_{ab}=1-\epsilon_n\), or \(\mathcal V_{\ell,q}=M_n\), including
an exceedance produced by a jump. Save the positive active credits
and freeze the process. If no guard is hit, \(\sigma=\infty\);
the deterministic terminal time \(T\) is not itself a failure event.

When a mixed observation is saved at a root's positive boundary, take
its value after geometric updates but before removing auxiliary roots
because of simultaneous positive retirement. At a continuous hit take
the boundary values before changing the flags. This convention prevents
a tag witness from disappearing at the very instant it is observed.
It does not alter the selected set or its birth rates. Auxiliary roots
that retired strictly earlier contribute zero.

\begin{proposition}[Construction and legality]\label{prop:construction}
For fixed \(n\), the stopped process is well-defined on \([0,T]\),
is adapted with right-continuous paths and left limits, and always
produces a no-three-in-line set of at most \(2n\) points.
\end{proposition}
\begin{proof}
Give each point a Poisson candidate clock of rate \(B/n\), together
with independent uniform marks. At a candidate arrival, accept the
birth if its mark is at most \(m_t(p)w_p(t-)/B\) and the point is
eligible in the pre-arrival state. The total candidate rate is \(nB\),
so there are finitely many arrivals on \([0,T]\) almost surely.
Use the filtration revealing only arrivals and marks up to the current
time, with the usual completion and right-continuous modification.

Between arrivals and mode changes, solve the finite-dimensional ODE.
Within the stopped region the denominators are separated from zero:
\(k\le K<2\), \(r\ge\rho>0\), and \(\eta\ge1-\epsilon_n>0\).
Using \(w\), rather than \(\log w\), as a coordinate, each mode has a
locally Lipschitz vector field. Extend it outside a compact
neighborhood with a cutoff if necessary. The resulting flow is unique
and depends continuously on its initial state. First hits of its
closed guards are measurable: on any bounded time interval, the
minimum distance of the continuous flow to a closed guard detects a
hit. Composition at the preceding stopping time preserves adaptation.

Every positive retirement occurs at most once per point. Occupancies
and line activation flags change only finitely often, deadlines and
phases are finite, and a global stop occurs once. Resolve all
simultaneous guards together and mark them as processed. Together with
the finite candidate arrivals, this rules out an accumulation of
updates. The construction gives the claimed path regularity.

Immediately before every birth the chosen point is active, hence is
not on a line with two selected points. Adding it preserves legality;
the subsequent deaths maintain this invariant. The row bound gives
\(|S|\le2n\).
\end{proof}

For compensators, evaluate the intensity in the pre-jump state and use
\(\ind_{t\le\sigma}\). The birth that triggers a global stop is counted.
In time integrals one may instead use \(\ind_{t<\sigma}\), since the
two differ on at most one time per path. Thinning gives, for a
nonnegative predictable observation \(g\),
\begin{equation}\label{eq:compensator}
 \E\sum_{\text{births }(p,t)}g_p(t)
 =\E\int_0^T\sum_p\lambda_p(t)g_p(t)\dd t.
\end{equation}
For simple predictable functions this is independent increments and
the uniform-mark rule; monotone approximation gives the general
identity. Bounded signed observations follow by subtraction. All
fixed-\(n\) rates and stopped observations used below are integrable.

\section{Saved credit and the exact output identity}\label{sec:loss}

Define \(Z_p\) to equal \(w_p\) while \(p\) is active, zero after
selection or geometric death, and its saved value after positive
retirement or a global stop. Let \(C_p\) be the birth-eligible credit,
including the global-stop indicator, and let \(V_p\) be saved credit.
Then \(Z_p=C_p+V_p\) and \(0\le Z_p\le B\).

\begin{lemma}[One-point compensation]\label{lem:credit}
For each point \(p\), \(Z_p\) is a bounded martingale and
\(\E Z_p(t)=1\).
\end{lemma}
\begin{proof}
While \(p\) is active, let \(J_p(q)\in\{0,1/2,1\}\) be its
multiplier at a birth of \(q\). Selection or deletion gives zero,
a first designated-line birth gives \(1/2\), and otherwise it is one.
The flow \eqref{eq:weightflow} is precisely
\[
 \dot w_p=w_p\sum_q\lambda_q(1-J_p(q)).
\]
By \eqref{eq:compensator}, the jump drift is its negative. Positive
saving changes no value; subsequent drift is zero. A coincident
geometric death has already been assigned multiplier zero. Thus the
compensated drift vanishes, and boundedness makes the result a true
martingale. Its initial value, including immediate saving of unused
points, is one.
\end{proof}

\begin{lemma}[Point law and birth intensities]\label{lem:pointlaw}
The formula \(\dd\Prob^{(p)}=Z_p(T)\dd\Prob\) defines a probability
measure. At any bounded stopping time \(\theta\le T\),
\(\E^{(p)}X=\E[Z_p(\theta)X]\) for nonnegative
\(\F_\theta\)-measurable \(X\). Before the root's saving boundary,
the birth intensity of \(q\) under this law is \(J_p(q)\lambda_q\).
\end{lemma}
\begin{proof}
The martingale identity gives
\(\E[Z_p(T)\mid\F_\theta]=Z_p(\theta)\).
One can obtain the stopping-time version from finite-valued dyadic
approximations from above, right continuity and bounded convergence.
For the intensity statement, enumerate the births, condition
\(Z_p(T)\) on each birth's post-jump sigma field, substitute
\(Z_p(t)=J_p(q)Z_p(t-)\), and then apply
\eqref{eq:compensator}. Tonelli justifies the nonnegative birth sum.
On time integrals \(Z_p(t-)=Z_p(t)\) almost everywhere.
\end{proof}

The changed law forbids a birth that kills the root, but does not
delete the self-rate term \(\lambda_p\) from the original deterministic
weight ODE. Losing this distinction changes the compensation identity.

\begin{proposition}[Exact remaining-exposure loss]\label{prop:loss}
With \(E_p\) and \(\theta_p\) as in Section~\ref{sec:overview},
\eqref{eq:previewloss} holds for every fixed \(n\). Therefore
\[
 \sup_p\Prob^{(p)}(E_p)\le\varepsilon
 \quad\Longrightarrow\quad
 \E|S_n|\ge(1-\varepsilon)n A_n.
\]
\end{proposition}
\begin{proof}
Apply \eqref{eq:compensator} to the birth count and use \(\E Z_p(t)=1\):
\[
 \E|S_n|=\frac1n\sum_p\int_0^T m_t(p)\E C_p(t)\dd t
 =nA_n-\frac1n\sum_p\int_0^T m_t(p)\E V_p(t)\dd t.
\]
On a path with positive saving,
\(V_p(t)=Z_p(T)\ind_{\theta_p\le t}\). Tonelli changes its integral
to \(Z_p(T)\int_{\theta_p}^T m_t(p)\dd t\).
Deadline saving has no remaining exposure. A path dying without
positive saving has zero terminal credit and may be assigned
\(\theta_p=T\). The density identity proves \eqref{eq:previewloss}.
Finally the remaining exposure is between zero and the deterministic
total \(F_{\typ(p),T}(p/n)\); sum the probability bound. Events need
not be independent or disjoint.
\end{proof}

\begin{lemma}[Product formula and a positive lower bound]\label{lem:product}
While \(p\) is active,
\begin{equation}\label{eq:product}
 w_p(t)=
 \exp\left\{\int_0^t H_p(u)\dd u+
 \int_0^t\lambda_p(u)\left(1-
       \sum_{\substack{\ell\ni p\\\ell\in\Gamma}}b_{N_\ell}\right)\dd u\right\}
 \prod_{\substack{\ell\ni p\\\ell\in\Gamma}}a_{N_\ell}(k_\ell).
\end{equation}
For large enough \(n\), \(w_p\ge2^{-g-1}\) uniformly while active.
The absolute value of the self-correction integral is at most
\(\max(1,g-1)TB/n=o(1)\).
\end{lemma}
\begin{proof}
Between births, \(\mathrm d\log a_N(k)/\mathrm dt=\Lambda_\ell b_N\).
At a first designated-line birth \(a_1/a_0=1/2\), exactly matching
the surviving point's halving. Subtract these factors from
\eqref{eq:weightflow}; all designated-line terms cancel except the
displayed self-correction. Initially every factor is one. Since
\(H_p\ge0\), \(a_0\ge1\), \(a_1\ge1/2\), and \(\lambda_p\le B/n\),
the asserted bounds follow.
\end{proof}


\section{Mixed moments without independence}\label{sec:moments}

The estimates must retain leading coefficient one. Replacing all
joint probabilities by an unspecified constant multiple of products
would lose the narrow numerical margin. The following compensating
factors separate ordinary-line dependence, designated-line dependence,
and simultaneous changes caused from outside a pair's baseline.

Fix disjoint sets of labels \(C,D\), called roots and witnesses, and
keep their union \(L=C\cup D\) fixed through the argument. If a witness
\(q\) has been born, let \(b_q\) be a deterministic factor
\(e_q(t_q)\in[0,1]\) evaluated at its birth time; otherwise set it to
zero. For a line put \(m_\ell=|C\cap\ell|\). Define
\begin{align}
 \Phi_{C,D}={}&\ind_{C\text{ active}}
 \prod_{q\in D}b_q\prod_{p\in C}w_p\,
 \prod_{\ell\notin\Gamma}r_\ell^{(m_\ell-1)_+}
 \prod_{\ell\in\Gamma}a_{N_\ell}(k_\ell)^{-(m_\ell-1)_+},
 \label{eq:phi}\\
 E_L={}&\prod_{\{a,b\}\subseteq L}\eta_{ab},
 \qquad Y_{C,D}=E_L\Phi_{C,D}.\label{eq:etaall}
\end{align}
Only lines containing two roots contribute nontrivial factors, so the
products are finite. The correction \(E_L\) uses the entire label set,
not just the current roots. This makes it identical on the two sides
of a witness-birth transition.

\begin{lemma}[Generator and source bound]\label{lem:source}
The positive source in the compensated drift of \(Y_{C,D}\) is at
most
\begin{equation}\label{eq:source}
 \sum_{q\in D}\beta_q(t)Y_{C\cup\{q\},D\setminus\{q\}},
\end{equation}
where \(\beta_q=m_t(q)e_q(t)/n\) if no pair of labels in \(L\) shares
a designated line. In general one may multiply each \(\beta_q\) by
\(\max(1,A_1)^g\), where \(A_1=1/(2-K)\).
Every other compensated drift and every auxiliary positive-retirement
jump is nonpositive.
\end{lemma}
\begin{proof}
Fix a birth cause outside the current roots and group roots by the
lines they share with it. On an active ordinary line, the \(m\) root
compensations contribute \(m\lambda_q\), while the factor
\(r^{m-1}\) subtracts \((m-1)\lambda_q\); the group is killed at the
birth. On a designated line of occupancy zero, the continuous
contribution is \(\lambda_q/2\), and the group multiplier is
\((1/2)^m(a_0/a_1)^{m-1}=1/2\). Occupancy one gives continuous
contribution \(\lambda_q\) and multiplier zero. An empty ordinary
line contributes nothing and has multiplier one.

Thus each group has a multiplier \(J_i\in\{0,1/2,1\}\), and its
combined relative drift is
\[
 \lambda_q\left[\sum_i(1-J_i)+\prod_iJ_i-1\right].
\]
For \(d_i=1-J_i\in[0,1]\), the finite Bonferroni inequality gives
\[
 0\le\sum_i d_i-\left(1-\prod_i(1-d_i)\right)
 \le\sum_{i<j}d_i d_j.
\]
One way to prove the upper bound is to introduce independent Bernoulli
variables with success probabilities \(d_i\) just for this finite
algebraic identity: if their sum is \(N\), then
\(N-\ind_{N\ge1}\le\binom N2\), and take expectations.
This auxiliary proof does not assume independence in the geometric
process. Different groups supply a pair of roots whose baseline does
not contain \(q\). The negative logarithmic drift of \(E_L\), namely
\(-\sum_{\{a,b\}\subseteq L}\kappa_{ab}\), pays this upper bound.
There is no jump cross term because \(E_L\) is continuous.
If \(q\) is a root, its own continuous compensation and zeroing jump
cancel; the shared-line correction cancels the other terms on its
line. Thus internal causes introduce no extra positive drift.

It remains to consider an unborn witness \(q\in D\). On a branch
where all roots survive its birth, any shared ordinary line was
previously empty, so its credit factor is one. A shared designated
line containing \(k\ge1\) roots must previously have occupancy zero.
Comparing the post-birth observation with the pre-birth source gives
the exact factor
\[
 2^{-k}a_1^{-(k-1)}a_0^k=a_1.
\]
The rate divided by the source's \(w_q\) is \(m_t(q)/n\).
The witness factor supplies \(e_q(t)\), and \(E_L\) cancels exactly.
There are at most \(g\) shared designated lines. This proves
\eqref{eq:source}. Prior retirement of an auxiliary root only drops
the observation to zero.
\end{proof}

\begin{proposition}[Triangular mixed-moment bound]\label{prop:mixed}
Write \(R=|L|\) and \(F_q^\beta(t)=\int_0^t\beta_q(u)\dd u\).
On the stopped process,
\begin{equation}\label{eq:mixed}
 \E\Phi_{C,D}(t)\le
 (1-\epsilon_n)^{-\binom R2}\prod_{q\in D}F_q^\beta(t).
\end{equation}
The same estimate holds when the entire fixed-label hierarchy is
frozen at one further bounded stopping time.
\end{proposition}
\begin{proof}
With \(D=\varnothing\), the initial value is at most one and there
is no positive source. With \(D\ne\varnothing\), the initial value
is zero. Set \(U_{C,D}=\E Y_{C,D}\). Lemma~\ref{lem:source} and the
compensator identity give
\[
 U_{C,D}(t)\le\sum_{q\in D}\int_0^t
 \beta_q(s)U_{C\cup\{q\},D\setminus\{q\}}(s)\dd s.
\]
Induct on \(|D|\). The resulting sum of integrals is exactly
\(\prod_{q\in D}F_q^\beta(t)\), by differentiating that finite
product and using its value zero at time zero. Finally
\(E_L\ge(1-\epsilon_n)^{\binom R2}\). At a common extra stop the
observations are saved, their subsequent sources vanish, and the
same nonnegative deterministic integrals remain upper bounds.
\end{proof}

For fixed \(R\), the first factor in \eqref{eq:mixed} is \(1+o(1)\).
Removing connection factors costs only a constant depending on
\(R,g,K,\rho\), because \(r\ge\rho\) and \(a_N\le2/(2-K)\).
In particular, products with \(j\) distinct born witnesses cost
at most \(C_R(t/n)^j\). A repeated root weight is bounded by
\(w^k\le B^{k-1}w\). A witness repeated many times is still just one
birth and must not be charged as independent births.

For estimates under \(\Prob^{(p)}\), put \(p\) among the roots of
every member of the hierarchy, and freeze all of it at the first
positive root boundary or global stop, capped at \(T\). Use the
postgeometry, prepositive observation convention. Lemma~\ref{lem:pointlaw}
then converts the expectation to the desired point-law moment.
No witnesses born after this common boundary may be added.
These conventions are needed for estimates at first crossings,
not just at deterministic times.


\section{Removing the global safety stops}\label{sec:global}

All estimates here are proved for the stopped process. On it,
\(k\le K\), \(r\ge\rho\), and \(\eta\ge1-\epsilon_n\), so
Proposition~\ref{prop:mixed} is available before any probability bound
on the global stop has been established.

\subsection{Pair-correction and virtual-deletion guards}
\begin{lemma}[Pair-correction tail]\label{lem:etatail}
For each fixed positive integer \(m\),
\begin{equation}\label{eq:etatail}
 \Prob(\text{pair-correction guard})
 \le C_m n^{4-m+2m\xi+d(m-1)+m\zeta}.
\end{equation}
\end{lemma}
\begin{proof}
For a fixed pair \(a,b\), classify a cause changing both roots
according to whether the two lines are ordinary or designated.
Two ordinary lines require two distinct selected witnesses; one
ordinary line requires one; two designated lines require only an
intersection of two among finitely many directions.

In an \(m\)-fold expansion, regard distinct selected witnesses as
vertices and the double-ordinary causes as edges of a multigraph.
If there are \(j\) witnesses and \(c\) connected components, the
first edge of a component costs \(C_\xi n^{2+\xi}\) by
\eqref{eq:twostar}. An edge adding one new witness costs
\(C_\xi n^{1+\xi}\) by \eqref{eq:fixedline}. With both witnesses
known, the cause is the intersection of two fixed lines and has at
most one choice. Coincidence of those lines would put the cause on
the excluded baseline. Thus the count is at most
\(C_m n^{j+(j-c)\xi}\le C_m n^{j+2m\xi}\).
A component containing a designated/ordinary cause may be started
with that cause, at cost \(C_\xi n^{1+\xi}\). A double-designated
cause has at most \(g^2\) choices. This also covers mixed components
and a witness appearing on different sides of different edges.

If the physical causes have \(k\ge1\) distinct labels, their rates
supply \(n^{-m}\), root repetitions cost
\(B^{m-k}\le B^{m-1}\), and the \(j\) distinct witnesses cost
\(C_m n^{-j}\) by Proposition~\ref{prop:mixed}. Root activity
indicators can be replaced by constant multiples of root weights
using Lemma~\ref{lem:product}. A label that is both an active cause
and a selected witness gives zero. Therefore
\[
 \E\kappa_{ab}(t)^m\le C_m n^{-m+2m\xi}B^{m-1}.
\]
Hitting the floor requires
\(\int_0^T\kappa_{ab}\dd t\ge-\log(1-\epsilon_n)\ge\epsilon_n\).
Apply Holder to the time integral, then Markov and a union bound
over at most \(n^4\) pairs. This is \eqref{eq:etatail}.
\end{proof}

\begin{lemma}[Virtual-deletion tail]\label{lem:virtualtail}
For each fixed positive integer \(m\),
\begin{equation}\label{eq:virtualtail}
 \Prob(\text{virtual-deletion guard})
 \le C_m\sum_{a=1}^m n^{3-\gamma m+d(m-a)+a\xi}.
\end{equation}
\end{lemma}
\begin{proof}
Fix a designated \(\ell\) and an active \(q\notin\ell\).
In an \(m\)-fold expansion of \eqref{eq:virtual}, suppose there are
\(a\) distinct leaves on \(\ell\). The lines from \(q\) to these
leaves are distinct, and their selected witnesses are distinct:
a common witness on two such lines would be \(q\), which is active.
By \eqref{eq:fixedline} the number of leaf/witness choices is at most
\(C_m n^{a(1+\xi)}\). Source prices supply \(n^{-a}\), and repeated
leaves cost \(B^{m-a}\). The normalized weight is at most \(2w\).
Consequently
\[
 \E[\ind_{q\text{ active}}\mathcal V_{\ell,q}^m]
 \le C_m\sum_{a=1}^m B^{m-a}n^{a\xi}.
\]
Apply the same fixed-label hierarchy at the first crossing; no
union over observation times is necessary. There are \(O(n)\)
designated lines and \(n^2\) possible external points. Markov at
level \(n^\gamma\) gives \eqref{eq:virtualtail}.
\end{proof}

The slope in \(m\) in \eqref{eq:etatail} is
\(-(1-d-2\xi-\zeta)<0\). Since \(\xi<d\), the largest exponent in
\eqref{eq:virtualtail} occurs at \(a=1\) and equals
\(3-d+\xi-m(\gamma-d)\). Both probabilities are therefore smaller
than any prescribed fixed power of \(n^{-1}\), after choosing one
sufficiently large fixed moment order. The process itself is unchanged
by that choice of proof order.

\subsection{The protected law of a designated line}
Fix a designated line and one of its occupancy stages \(N=0,1\).
Starting at the stage's entry time \(U\), suppress its next birth
through the density process
\[
 D_t=\ind_{\text{no new line birth in }(U,t]}
 \exp\left(\int_U^t\Lambda_\ell(s)\dd s\right),
\]
stopped and saved at the relevant guards or \(T\). Before \(U\) its
value is one. The birth at \(U\), if any, starts the stage and is
not counted again as its next birth. The continuous drift
\(D\Lambda_\ell\) cancels the zeroing-jump drift. Also
\[
 \int_U^t\Lambda_\ell\dd s
 =\int_{k(U)}^{k(t)}a_N(k)\dd k
 \le2\log\frac2{2-K}.
\]
Thus \(D\) is a bounded martingale and defines a probability law
conditionally on the entry history. Its upper bound is the fixed
constant \(A_0^2\), where \(A_0=2/(2-K)\). On a crossing occurring
before the next line birth, \(D\ge1\), so the original crossing
probability is at most its probability under the protected law.

Normalize each active line weight by
\(c_p=w_p/a_N(k_\ell)\), and save this \emph{normalized} value at
positive retirement. It is not subsequently divided by a changing
\(a_N\). For a fixed cell/type subset of the line, let \(M\) be the
sum of current and saved \(c_p\). Define its entropy sum using
\[
 h(c)=c\left(1+\log\frac{2B}{c}\right),\quad h(0)=0,
 \qquad H=\sum h(c_p),
\]
again with saved terms frozen. Under the protected law the surviving
positive drift of each credit is
\(g_p=\lambda_p(1-b_N)c_p\). This follows by subtracting
\(\mathrm d\log a_N/\mathrm dt=\Lambda_\ell b_N\) from the original
ODE and deleting only the line-birth jump rates. It is
\(\lambda_pc_p/2\) at occupancy zero and zero at occupancy one.

\begin{lemma}[Credit, entropy and whole-event energy]\label{lem:entropy}
Stop further at \(M\le C_0n\). There are compensated martingales
\(Y_M,Y_H\) and a nonnegative cumulative cost \(V\) such that
\begin{align}
 M&=M(U)+D_M+Y_M,&D_M&=O(B),\label{eq:creditdecomp}\\
 H&=H(U)+D_H-V+Y_H,&D_H&=O(B\log n),\label{eq:entropydecomp}
\end{align}
and, before the virtual guard,
\begin{align}
 \langle Y_M\rangle&\le C n^\gamma V,
 &|\Delta Y_M|&\le C n^\gamma,\label{eq:bracketM}\\
 \langle Y_H\rangle&\le C n^\gamma(\log n)^2V,
 &|\Delta Y_H|&\le C n^\gamma\log n.\label{eq:bracketH}
\end{align}
The constants are independent of the admissible entry history.
\end{lemma}
\begin{proof}
For a multiplier \(J\in\{0,1/2,1\}\), direct algebra gives
\[
 h'(c)c(1-J)+h(Jc)-h(c)=-c(J\log J-J+1).
\]
The cost is \(c\) for deletion, \(c(1-\log2)/2\) for halving,
and zero otherwise. It is comparable by fixed positive constants
to the credit decrease \(c(1-J)\). Define \(V\) to be the rate
integral of the \emph{sum of these costs for each complete birth
event}. Compensating gives \eqref{eq:creditdecomp} and
\eqref{eq:entropydecomp}. Since \(\lambda_p\le B/n\) and the
credit sum is capped by \(C_0n\), the first drift is \(O(B)\).
The active credits have the positive lower bound supplied by
Lemma~\ref{lem:product} divided by \(A_0\), so
\(h'(c)=\log(2B/c)=O(\log n)\) on them. This gives the entropy
drift estimate.

In one external birth event, the ordinary-line deletion of normalized
mass is at most \(M_n\), and effects through other designated
directions touch at most \(g\) points, each of credit at most \(2B\).
Its total credit decrease \(\Delta\) is therefore at most
\(M_n+2gB=O(n^\gamma)\). The square of the \emph{whole} decrease
satisfies \(\Delta^2\le C n^\gamma\Delta\).
The absolute entropy decrease is at most \(C\log n\) times this
credit decrease, including deletion to zero. Comparing the decrease
to the event's total cost and integrating its birth rate proves
\eqref{eq:bracketM}--\eqref{eq:bracketH}. These estimates retain
all cross terms of simultaneous changes. Summing the squares of
individual point changes would not prove the required bracket bound.
\end{proof}

\begin{lemma}[Exponential crossing estimate]\label{lem:exp}
Let \(Y\) be a compensated birth martingale with \(Y_0=0\), jumps
bounded in absolute value by \(b\), and continuous predictable
quadratic variation \(Q\). For \(x,v>0\),
\begin{equation}\label{eq:freedman}
 \Prob(\exists t\le T:Y_t\ge x,\ Q_t\le v)
 \le\exp\left(-\frac{x^2}{2(v+bx/3)}\right).
\end{equation}
\end{lemma}
\begin{proof}
For \(0<\theta b<3\) and \(|u|\le b\), the Taylor series and
\(k!\ge2\cdot3^{k-2}\) for \(k\ge2\) give
\[
 e^{\theta u}-1-\theta u
 \le\psi u^2,\qquad \psi=\frac{\theta^2}{2(1-\theta b/3)}.
\]
Stop at the first \(Y\ge x\), capped at \(T\). The stopped
exponential \(\exp(\theta Y-\psi Q)\) is bounded by
\(\exp(\theta(x+b))\). Its compensated drift is nonpositive by the
displayed inequality, so its expectation is at most one.
Since \(Q\) is continuous and nondecreasing, a crossing with
\(Q\le v\) exists precisely when the first crossing has \(Q\le v\).
On this event the terminal stopped exponential is at least
\(\exp(\theta x-\psi v)\). Markov gives
\(\exp(-\theta x+\psi v)\); take \(\theta=x/(v+bx/3)\).
This proof uses a bounded stopped exponential, rather than assuming
that an unlocalized exponential local martingale has mean one.
\end{proof}

\begin{proposition}[Clock stops are rare]\label{prop:clock}
For each fixed positive tolerance and admissible fixed data, the
failure probability of the simultaneous designated-line credit
bound is at most
\[
 Cn\exp\left(-c\frac{n^{1-\gamma}}{\log n}\right).
\]
In particular, the probability of a designated clock reaching \(K\)
is smaller than every fixed inverse power of \(n\).
\end{proposition}
\begin{proof}
Stop the cumulative cost at \(C_1n\log n\). Since \(H\ge0\),
\(H(U)=O(n\log n)\), and \(D_H=o(n\log n)\), reaching this cost
cap, for sufficiently large \(C_1\), forces an upward fluctuation
of \(Y_H\) of order \(n\log n\). Apply Lemma~\ref{lem:exp} with
\(v=O(n^{1+\gamma}(\log n)^3)\) and
\(b=O(n^\gamma\log n)\). An increase of \(M\) by a fixed multiple
of \(n\) similarly forces a fluctuation of \(Y_M\) of order \(n\),
with \(v=O(n^{1+\gamma}\log n)\) and \(b=O(n^\gamma)\).
Both exponents have order \(n^{1-\gamma}/\log n\).

At the first birth on the line,
\((w/2)/a_1=w/a_0\); the selected point is removed. The saved ledger
is unchanged and the moving normalized credit cannot increase.
At the second birth, all remaining moving credit disappears. To
treat the occupancy-one stage, condition on its actual entry history
and the preceding good credit bound. The entropy bound follows
from that credit bound and \(c\le2B\). The same constants therefore
apply to every such history. Average the conditional estimate; do
not take a union over possible histories or starting times.

There are \(O(n)\) designated lines and a fixed number of cell/type
pieces and phases. Summing their crossing probabilities preserves
the displayed bound. On the resulting good event, integrate the
moving normalized credit with the masks to obtain
\[
 k_\ell\le\text{deterministic line capacity}
             +C_{\mathrm{data}}T\varepsilon+O(n^{-1}).
\]
Choose \(\varepsilon>0\) first so that the second term is smaller
than half of \(\kappa_H\), then choose \(n\) large. The capacity
condition \eqref{eq:heavyinput} precludes hitting \(K\).
Positive retirement need not be rare for this argument: its credit
was retained in the ledger throughout.
\end{proof}

Combining the three guards shows, for every fixed \(A_0'>0\),
\begin{equation}\label{eq:globaltail}
 \Prob(\sigma\le T)\le C_{A_0'}n^{-A_0'}.
\end{equation}
Under any point law this is at most \(B\) times as large, hence is
\(o(1)\) uniformly in the root. This establishes global-stop smallness
before the weight-cap and ordinary-line tag estimates below.




\section{Ordinary-line risk under the point law}\label{sec:riskproof}

Fix a root \(p\) and its deterministic deadline \(\tau=\tau_p\).
All observations below are stopped and saved at the root's first
positive boundary or global stop. A quantity explicitly restricted
to the active root is zero afterwards. Write
\(F_q(t)=\int_0^t m_u(q)\dd u\), and let
\[
 I_p=\int_0^{\theta_p\wedge T}\ind_{p\text{ active}}H_p(t)\dd t
\]
be its integrated ordinary-line hazard. Earlier root death has
\(\Prob^{(p)}\)-probability zero. We will bound \(I_p\) by a proxy
having sharp fixed-order moments, plus an error small in \(L^1\).
The distinction between those two kinds of estimate is essential.

\subsection{Removing ordinary-line credit at additive cost}
For an ordinary line through \(p\), let
\(H_\ell(t)=r_\ell(t)\sum_{q\in\ell\setminus\{p\},\ q\text{ active}}w_q(t)\).
The two-root, no-witness case of Proposition~\ref{prop:mixed}, followed
by the conditional density identity, gives
\begin{equation}\label{eq:conditionalH}
 \E^{(p)}[\ind_{p\text{ before boundary at }t}H_\ell(t)\mid\F_u]
 \le(1+o(1))\ind_{p\text{ before boundary at }u}H_\ell(u)
 \quad(u\le t).
\end{equation}
The correction ratio is bounded by a fixed power of
\((1-\epsilon_n)^{-1}\). Conditional division is performed only
where \(Z_p(u)>0\); its complement is null under the point law.
The auxiliary-root sum excludes \(p\), so no self-square is
silently included in this two-root argument.

One selected witness \(z\in\ell\) and two leaves give
\begin{equation}\label{eq:secondH}
 \E^{(p)}[\ind_{p\text{ before boundary},\ z\in S_t}H_\ell(t)^2]
 \le C L_\ell(L_\ell+B)t/n.
\end{equation}
For distinct leaves the connection factor is exactly \(r_\ell^2\);
a repeated leaf is paid using \(r^2w^2\le Brw\).
The one-leaf version is at most \(CL_\ell t/n\).
All labels here lie on the same ordinary line, so there is no
designated-line source multiplier.

Since \(r\) starts moving at the first line birth,
\(1-r_\ell(t)=\int_{\text{activation}}^t r_\ell\Lambda_\ell\dd u\).
The correct self-rate bound is
\begin{equation}\label{eq:selfrate}
 r_\ell\Lambda_\ell\le H_\ell/n+B/n.
\end{equation}
Multiplying the leaf sums in the hazard by \(r_\ell\) therefore loses
a nonnegative amount whose expectation is at most
\begin{equation}\label{eq:removecredit}
 Cn^{-3}\sum_{\ell\ni p}(L_\ell^3+BL_\ell^2)
 +CBn^{-3}\sum_{\ell\ni p}L_\ell^2
 =O(1+(B/n)\log n)=o(\log n).
\end{equation}
For completeness, insert the integral for \(1-r\), use
\eqref{eq:conditionalH} to move the future \(H\) to the earlier
time, cover activation by \(\sum_{z\in\ell}\ind_{z\in S_u}\),
and apply \eqref{eq:secondH}. The second term in
\eqref{eq:removecredit} is the separate self-rate contribution from
\eqref{eq:selfrate}. The line sums are Lemma~\ref{lem:lines}.
This is an additive error, not a constant-factor loss in the leading
\(\log n\) coefficient.

\subsection{Long lines: replace the leaf sum, retain the birth time}
Call a line long when \(L_\ell\ge B\log n\). Fix a selected witness
\(z\in\ell\), and put
\[
 H=r_\ell\sum_{q\in\ell\setminus\{p\},\ q\text{ active}}m_t(q)w_q,
 \qquad D_z(t)=\sum_{q\in\ell\setminus\{p,z\}}m_t(q).
\]
For every fixed integer \(k\), the mixed hierarchy gives
\begin{equation}\label{eq:longmoment}
 \E^{(p)}[\ind_{p\text{ before boundary},\ z\in S_t}H^k]
 \le(1+o(1))\frac{F_z(t)}n
       [D_z(t)^k+C_kB L_\ell^{k-1}].
\end{equation}
Indeed, the all-distinct leaves have exactly the required connection
\(r_\ell^k\) and unit source price. A term with \(a<k\) distinct
leaves costs \(B^{k-a}L_\ell^a\le BL_\ell^{k-1}\), since
\(B\le L_\ell\). There are finitely many repetition patterns for
fixed \(k\). No division by a positive lower bound on \(F_z\) is
used; if it is zero the corresponding birth event is null.

For fixed \(\epsilon>0\), use
\(H\ind_{H>D+\epsilon L}\le H^k/(D+\epsilon L)^{k-1}\).
Since \(D\le L\), the main term is at most
\(L(1+\epsilon)^{-(k-1)}\), while the repeated-leaf term is
\(C_{k,\epsilon}B\). Sum over witnesses, long lines and time,
with the hazard factor \(1/n\). As \(B/L\le1/\log n\), the
positive error \(R_{\rm long}\) made by replacing the leaf sum by
\(D_z(t)\) satisfies
\begin{equation}\label{eq:longerror}
 \limsup_{n\to\infty}\sup_p
 \frac{\E^{(p)}R_{\rm long}}{\log n}
 \le C\epsilon+C(1+\epsilon)^{-(k-1)}.
\end{equation}
First choose \(\epsilon\), next a fixed sufficiently large \(k\),
and only then let \(n\) grow. This proves an \(o(\log n)\)
\(L^1\) error for the same process.

For a witness born at time \(u\), define its deterministic remaining
cost and birth price by
\[
 A_z(u)=\frac1n\int_u^\tau D_z(t)\dd t\le\frac{TL_{pz}}n,
 \qquad b_z=\frac1n\int_0^\tau m_u(z)A_z(u)\dd u.
\]
Let \(P_{\rm long}\) be the sum of \(A_z\) over long-line witnesses
born before the root boundary, and let \(M_{\rm long}=\sum_z b_z\).
Insert \(A_z(u)/(TL_{pz}/n)\le1\) as a deterministic birth-retention
factor in the mixed hierarchy. It follows that for every fixed \(m\),
\begin{equation}\label{eq:longproxy}
 \|P_{\rm long}\|_{L^m(\Prob^{(p)})}
 \le M_{\rm long}+o(\log n).
\end{equation}
Here are the non-generic terms in this assertion. A repeated witness
block of size \(r\ge2\) has total price at most
\(C_r n^{-(r+1)}\sum_\ell L_\ell^{r+1}=O_r(1)\), whereas a
singleton block costs \(O(\log n)\). Thus repetitions lose a power
of \(\log n\). For two distinct witnesses sharing a designated
line, that line cannot pass through the root. By
\eqref{eq:fixedline}, the weighted configuration count is
\[
 \sum_{z_1}L_{pz_1}
 \sum_{\substack{z_2\text{ on a designated}\\\text{line through }z_1}}L_{pz_2}
 \le C n^{3+\xi}\log n.
\]
The two birth-and-cost factors supply \(n^{-4}\), giving
\(O(n^{-1+\xi}\log n)\). Remaining singleton factors cost
\(O(\log n)\) each. A finite union over pairs covers multiple
exceptions. The all-distinct, unshared case keeps the exact product
of the \(b_z\), with leading factor \(1+o(1)\).

\subsection{Short stars and their exceptional configurations}
An active short line through \(p\) has one selected witness \(z\).
Each active leaf \(q\ne p\) on that line gives a star \((p,q,z)\).
Let
\[
 X_t=\sum_{\text{short stars}}m_t(q)r_{\ell(p,q)}w_q,
 \qquad
 D_t=\frac1n\sum_{\text{deterministic short stars}}m_t(q)F_z(t).
\]
The actual sum is zero after the root boundary. In the deterministic
sum, \(p,q,z\) are distinct and collinear on a short ordinary line;
there is no selected-point condition.
Put \(t_0=n^{\gamma-1}\).

\begin{lemma}[Short-star moment estimate]\label{lem:short}
For each fixed positive integer \(m\), uniformly for \(t_0\le t\le T\)
and all roots,
\begin{equation}\label{eq:shortmoment}
 \E^{(p)}X_t^m\le(1+o(1))D_t^m+\epsilon_m(n)(tn\log n)^m,
 \qquad \epsilon_m(n)\longrightarrow0.
\end{equation}
\end{lemma}
\begin{proof}
In an \(m\)-fold product let \(a\) be the number of distinct leaves
and \(j\) the number of distinct witnesses. The active root makes
the selected witness on each of its lines unique, so
\(1\le j\le a\le m\). A leaf coinciding with a witness or the root
gives zero. If \(j<m\), the first leaf and witness of each star
line have \(O(n^2\log n)\) choices by \eqref{eq:line2}.
Additional leaves on that line have at most \(B\log n\) choices,
and repeated weights cost \(B^{m-a}\). The source price is
\(C_m(t/n)^j\). Dividing this bound by \((tn\log n)^m\) gives
\[
 C_m\left(\frac B{nt}\right)^{m-j}(\log n)^{a-m}=o(1),
\]
uniformly since \(B/(nt)\le n^{d-\gamma}\).

It remains to treat \(j=a=m\), with distinct star lines. The
following list exhausts the extra connections among different stars.
\begin{enumerate}[label=(\roman*)]
\item Two labels share a designated line. Fix the first label.
The line is outside the root, so \eqref{eq:fixedline} bounds the
second label together with its star partner by \(C n^{1+\xi}\),
instead of the generic \(O(n^2\log n)\) row. This covers leaf--leaf,
leaf--witness and witness--witness sharing.
\item A line connecting two leaves has length greater than
\(n^\gamma\). For a fixed first leaf there are
\(O(n^{2-\gamma})\) possible second leaves: sum the height counts
up to \(O(n^{1-\gamma})\). Its star witness has at most
\(B\log n\) choices. Relative to a generic star row the saving
is \(O(n^{d-\gamma})\).
\item A leaf--leaf line of length at most \(n^\gamma\) is active
through a new witness. There are at most \(n^\gamma\) positions
for it and an additional source price \(C_mt/n\), giving
\(O(n^{\gamma-1})\).
\item That line is active through an existing third star witness.
It cannot be either leaf's own witness, since then their star lines
would coincide. For a third witness use
\(\sum_{z\in\ell(q_i,q_j)}L_{pz}\le C n^{1+\xi}\), keeping all
\(m\) original birth prices. This again saves \(n^{-1+\xi}\).
An active leaf or the root cannot itself be a selected witness.
\end{enumerate}
Outside these cases the additional ordinary-line credits are all
one, there are no designated-line connections, and the star credits
exactly match the mixed-moment factors. The source price is the
exact product of \(F_z(t)/n\); keeping the masks produces \(D_t^m\).
Each single exception is bounded after dropping other restrictions,
so a finite union over label pairs also covers overlapping exceptions.
Removing remaining factors costs only \(C_m\), including
\(\rho^{-O(m^2)}\). The savings above and the correction
\((1-\epsilon_n)^{-O(m^2)}=1+o(1)\) prove the assertion.
\end{proof}

Minkowski's inequality now yields
\begin{equation}\label{eq:shortintegral}
 \left\|\int_{t_0}^\tau X_t/n\dd t\right\|_m
 \le\int_{t_0}^\tau D_t/n\dd t+o(\log n),
\end{equation}
with the integral zero when \(\tau<t_0\).
For earlier times the one-witness, one-leaf bound gives
\(\E^{(p)}X_t=O(tn\log n)\). Its integrated contribution is
\(O(t_0^2\log n)=o(\log n)\) in \(L^1\). No higher moment for
this early-time error is asserted.

\subsection{Hazard concentration and the weight cap}
Collect the credit-removal error, long-line leaf replacement error,
and early short-star contribution into \(R_n\ge0\), so
\(\sup_p\E^{(p)}R_n=o(\log n)\). Define
\[
 Y_n=P_{\rm long}+\int_{t_0}^\tau X_t/n\dd t.
\]
Then \(I_p\le Y_n+R_n\), and \eqref{eq:longproxy} and
\eqref{eq:shortintegral} give, for every fixed \(m\),
\[
 \|Y_n\|_m\le M_{\rm long}+M_{\rm short}+o(\log n)
 \le\mathcal B_n^{\Gamma^c}(p;F_\tau)+o(\log n).
\]
The last inequality is \eqref{eq:timesym}: the division into long
and short lines is symmetric in leaf and witness, and all four
fields use the same \(\tau\).
Proposition~\ref{prop:pair} and \eqref{eq:riskinput} therefore give
\(\|Y_n\|_m\le(r_*+o(1))\log n\), uniformly in the root.

For any \(b>r_*\), choose \(\epsilon>0\) with
\(r_*+\epsilon<b-\epsilon\). Markov on \(R_n\), followed by
the fixed-\(m\) moment bound, shows
\[
 \limsup_{n\to\infty}\sup_p\Prob^{(p)}(I_p>b\log n)
 \le\left(\frac{r_*+\epsilon}{b-\epsilon}\right)^m.
\]
The right side tends to zero as the fixed proof order \(m\) is
chosen larger. This quantifier order does not require constants
uniform over all \(m\), or an order depending on \(n\).
By \eqref{eq:product}, hitting \(B=n^d\) requires
\(I_p\ge d\log n-O(1)\). Choose \(r_*<b<d\). We have proved
\begin{equation}\label{eq:captail}
 \sup_p\Prob^{(p)}(\text{weight cap is an early saving cause})=o(1).
\end{equation}


\section{Ordinary-line tags and the final loss estimate}\label{sec:closing}

The ordinary-line tag \(r_\ell=\rho\) must be controlled uniformly
over all directions. A crude multiplication of a per-line bound by
\(n^2\) would lose too much. Instead retain the line's height and
count only the spatial cells that the particular line can meet.

\begin{lemma}[Linewise cell count]\label{lem:cells}
A nonaxis line meets at most \(2G-1\) open cells of a fixed
\(G\times G\) grid. Its lattice points have at most two parity types.
Consequently there are at most \(C_{\rm line}=4G\) nonempty cell/type
pieces on it, uniformly in \(n\), root and direction.
\end{lemma}
\begin{proof}
Inside the open square a nonaxis line crosses each of the
\(2(G-1)\) internal grid lines at most once. These crossings divide
it into at most \(2G-1\) intervals. Convexity prevents re-entry
into a cell. Coincident crossings at a vertex only reduce this count.
The integer points are \(p+kv\), with types alternating between
\(r\) and \(r+\bar v\). Thus there are at most
\(2(2G-1)\le4G\) pieces. Their index set is determined before the
random process is run.
\end{proof}

For an ordinary line \(\ell\ni p\) and one such piece \(E\), put
\[
 H_{\ell,E}(t)=r_\ell(t)
       \sum_{\substack{q\in(\ell\cap E)\setminus\{p\}\\q\text{ active}}}w_q(t),
 \qquad L_{\ell,E}=| (\ell\cap E)\setminus\{p\}|.
\]
Stop the fixed-label hierarchy at the first crossing
\(H_{\ell,E}>L_{\ell,E}+en\), or the root boundary, whichever
comes first. As usual, observations at a boundary precede changes
of positive-retirement flags. For fixed \(m\), let
\(\Theta_m(n)=(1-\epsilon_n)^{-\binom{m+1}2}\).
Expansion by the number \(a\) of distinct leaves gives
\begin{equation}\label{eq:tagmoment}
 \E^{(p)}H_{\ell,E}^m
 \le\Theta_m(n)\sum_{a=1}^m
 \left\{\begin{matrix}m\\a\end{matrix}\right\}
 B^{m-a}L_{\ell,E}^a.
\end{equation}
The braces are Stirling numbers of the second kind. For distinct
leaves the exact connection is \(r^a\); for repetitions use
\(r^m\le r^a\), \(w^j\le B^{j-1}w\). There are no witnesses and
no designated-line connections, so no extra fixed multiplier is
needed. Ordered distinct-leaf tuples are counted by
\((L)_a\le L^a\).

\begin{lemma}[Finite height sum]\label{lem:height}
For \(e>0\), integers \(m\ge3\) and \(N\ge0\),
\begin{equation}\label{eq:height}
 \sum_{h=1}^N4h\left(\frac2{2+eh}\right)^m
 \le\frac{16}{e^2}\left(\frac2{2+e}\right)^{m-3}.
\end{equation}
\end{lemma}
\begin{proof}
Put \(a=e/2\). For \(h\ge1\),
\[
 h(1+ah)^{-m}\le\frac{(1+a)^{-(m-3)}}{a(1+ah)^2},
\]
and
\[
 (1+ah)^{-2}\le\frac1a
 \left[\frac1{1+a(h-1)}-\frac1{1+ah}\right].
\]
Sum the latter telescoping difference. Multiplication by four
gives \eqref{eq:height}, since \(4/a^2=16/e^2\).
\end{proof}

Apply Markov to \eqref{eq:tagmoment}, sum first over the at most
\(C_{\rm line}\) pieces of each line, then over directions of
height \(h\), using \(4h\) directions and \(L\le2n/h\).
The all-distinct term is bounded by Lemma~\ref{lem:height}.
Bounding the denominator in each repeated term by \((en)^m\)
gives the explicit remainder
\begin{align}\label{eq:tagremainder}
 R_m(n;e)=4e^{-m}\bigg[&2\left\{\begin{matrix}m\\1\end{matrix}\right\}
                 n(B/n)^{m-1}
 +4\left\{\begin{matrix}m\\2\end{matrix}\right\}
                 (1+\log n)(B/n)^{m-2}\notag\\
 &+\sum_{a=3}^{m-1}2^a
 \left\{\begin{matrix}m\\a\end{matrix}\right\}
 \left(1+\frac1{a-2}\right)(B/n)^{m-a}\bigg].
\end{align}
Thus the probability of any such first crossing, under the fixed
root's point law, is at most
\begin{equation}\label{eq:tagprob}
 C_{\rm line}\Theta_m(n)
 \left[\frac{16}{e^2}\left(\frac2{2+e}\right)^{m-3}
       +R_m(n;e)\right].
\end{equation}
No union over times, or over all roots, has been used.
For sufficiently large fixed \(m\), \(R_m(n;e)\to0\); it is enough
that \(1-(1-d)(m-1)<0\) and \(m>2\).

On the complementary good event, sum the cell bounds after multiplying
by their deterministic mask values, then integrate \eqref{eq:tagflow}.
The self point costs a further \(TB/n\). The result is
\begin{equation}\label{eq:tagcapacity}
 1-r_\ell\le\text{deterministic ordinary-line capacity}
       +C_{\rm line}Te+TB/n+O(n^{-1}).
\end{equation}
Choose \(e=(1-\rho-\lambda)/(4TC_{\rm line})\) when \(T>0\).
The zero-time case has no births and is trivial. For large \(n\),
the last two terms are at most \((1-\rho-\lambda)/4\).
Using capacity at most \(\lambda\), the right side is then strictly
less than \(1-\rho\), precluding a tag. For any requested positive
error, choose \(m\) large enough that the first term of
\eqref{eq:tagprob} is below half that error, and then \(n\) large
enough for \(\Theta_m(n)\) and the remainder. This proves
\begin{equation}\label{eq:tagtail}
 \sup_p\Prob^{(p)}(\text{tag is an early saving cause})=o(1).
\end{equation}
The linewise count is used only here; it does not replace the number
of cells in global BV estimates or in the designated-line argument.

\begin{proof}[Completion of the analytic reduction]
The weight-cap estimate \eqref{eq:captail}, tag estimate
\eqref{eq:tagtail}, and global-stop estimate \eqref{eq:globaltail}
give \(\sup_p\Prob^{(p)}(E_p)=o(1)\). Their events may overlap;
a union bound suffices. Proposition~\ref{prop:loss} implies
\(\E|S_n|\ge(1-o(1))nA_n\), whereas the exact nonnegative loss
identity also gives \(\E|S_n|\le nA_n\). By Lemma~\ref{lem:chord},
\(A_n=A+O(n^{-1})\), proving Theorem~\ref{thm:bridge}.

For fixed \(c<A\) the expectation exceeds \(cn\) for all large
\(n\). If every possible legal output had cardinality smaller than
\(cn\), its expectation would be smaller too, a contradiction.
Equivalently, with \(k=\lceil cn\rceil\) and \(|S_n|\le2n\),
\[
 \Prob(|S_n|\ge k)\ge
 \frac{\E|S_n|-(k-1)}{2n-(k-1)}>0.
\]
This is an existence argument; it does not supply a practical
algorithm producing the target density on moderate grids.
\end{proof}




\section{From finite rectangles to certified geometric bounds}\label{sec:certificates}

The analytic reduction needs uniform inequalities on entire root
regions and on all line directions. A numerical value attached to a
sampled root is insufficient. This section describes a finite
rational certificate contract and proves its mathematical soundness.
Whether every concrete certificate has been supplied and checked is
a separate question, addressed in Section~\ref{sec:status}.

\subsection{A uniform chord envelope on a closed root box}
Let \(F\) be a nonnegative rational rectangular step field bounded by
\(M\), with zero outside the unit square. For direction \((1,u)\),
\(u>0\), write
\[
 C(v;u)=\int_\R F(t,v+ut)\dd t.
\]
For a closed root box \(Q_0=[a_0,a_1]\times[b_0,b_1]\), every root
has intercept in
\(l(u)=b_0-u a_1\le v\le h(u)=b_1-u a_0\).
For now assume \(Q_0\) is not a single point.

At each grid vertex \((a,\eta)\), with exterior cells read as zero,
define the signed corner coefficient
\[
 k_{a,\eta}=F_{\rm LU}+F_{\rm RD}-F_{\rm RU}-F_{\rm LD}.
\]
\begin{lemma}[Corner representation]\label{lem:corners}
For \(u>0\),
\begin{equation}\label{eq:corners}
 C(v;u)=\frac1u\sum_{a,\eta} k_{a,\eta}(v-\eta+ua)_+.
\end{equation}
\end{lemma}
\begin{proof}
For the indicator of \((a,b)\times(c,d)\), integrate
\(\ind_{v+ut>c}-\ind_{v+ut>d}\) over \(a<t<b\).
The first integral is
\([(v+ub-c)_+-(v+ua-c)_+]/u\), with the analogous expression
for \(d\). The four corner coefficients are \(-1,+1,+1,-1\).
Add these identities with the rational cell weights and combine
shared vertices. Nonaxis chords meet the zeroed grid boundaries
in finitely many points, so those values do not change the integral.
\end{proof}

For a vertex with negative coefficient set
\[
 P=a_1-a+\frac{\eta-b_0}{u},\qquad
 Q=a-a_0+\frac{b_1-\eta}{u},\qquad
 T_{a,\eta}(u)=\max\left(0,\min\left(P,Q,\frac{P+Q}{4}\right)\right).
\]
Here \(Q\) is a scalar; the root box is \(Q_0\).
Define
\begin{align}
 U_1(u)&=\max(C(l;u),C(h;u))+
             \sum_{k_{a,\eta}<0}(-k_{a,\eta})T_{a,\eta}(u),
             \label{eq:U1}\\
 U_2(u)&=C(l;u)+\sum_{J>0}J\,
       \big|[a,b]\cap[(\eta-h)/u,(\eta-l)/u]\big|.
       \label{eq:U2}
\end{align}
In \eqref{eq:U2} the sum runs over horizontal grid edges
\((a,b)\times\{\eta\}\) with upward jump
\(J=F_{\rm up}-F_{\rm down}>0\).

\begin{proposition}[Closed-box envelope]\label{prop:envelope}
For every root \((x,y)\in Q_0\), \(v=y-ux\), and \(u>0\),
\begin{equation}\label{eq:envelope}
 0\le C(v;u)\le U(u):=\min(U_1(u),U_2(u),M)\le M.
\end{equation}
The envelope is a finite expression formed from rational functions
\(\alpha+\beta/u\), rational arithmetic, minimum, maximum and
positive part.
\end{proposition}
\begin{proof}
For a corner term \(k(v-v_0)_+/u\), compare it with its endpoint
linear interpolation on \([l,h]\). A nonnegative \(k\) has no
positive excess. A negative \(k\) has excess zero if the corner
lies outside the interval, and otherwise maximum
\[
 (-k)\frac{(v_0-l)(h-v_0)}{u(h-l)}
 =(-k)\frac{PQ}{P+Q}
 \le(-k)\min(P,Q,(P+Q)/4).
\]
Sum the corner identities. Interpolation commutes with finite
addition, and the interpolated total is at most its larger endpoint.
This gives \(U_1\).

For fixed horizontal coordinate \(t\), move vertically from
\(l+ut\) to \(v+ut\). The field difference is the sum of its
crossed horizontal jumps. Discard negative jumps and extend the
upper endpoint to \(h+ut\). Integrating the resulting indicator
inequality over \(t\) gives exactly \eqref{eq:U2}. Boundary
exceptions involve finitely many values of \(t\). The trivial
bound \(C\le M\) follows because the major-coordinate length
inside the square is at most one. Each of the three bounds is
nonnegative. The displayed expressions for the intersections and
corners have the asserted finite form.
\end{proof}

For a point root, simply use \(C(l;u)\). A root box of zero width
or zero height is covered by the same argument whenever \(h-l>0\).
At \(u=0\) one may set \(U(0)=M\) for the angular integral.
This does not remove the separate obligation to check axis capacities.

\subsection{Verifying branches without trusting the branch generator}
Suppose a proposed branch on a rational interval \([r,s]\subset(0,1]\)
is \(V(u)=a+b/u\). To compare it with
\(a'+b'/u\), multiply the difference by positive \(u\).
Its sign is the sign of the affine function
\((a-a')u+(b-b')\), and is therefore determined by its two endpoint
values. Induction on the finite expression tree proves every min/max
branch choice throughout the interval. Check that the proposed
intervals cover \([0,1]\) in order, with no gaps in their interiors.

A finite exact branch partition exists: on each common child
interval, two nonidentical branches cross at at most one rational
point. Add that point if it lies inside the interval; identical
branches require no subdivision. Near zero, the final bounded
nonnegative envelope must have \(b=0\) on its first branch. A checker
can either verify this or use the constant bound \(M\) on that first
interval. It need not trust the search procedure that supplied the
partition.

\subsection{Rational integral bounds without logarithms}
On \(0<r<s\), suppose \(V(u)=a+b/u\) has nonnegative endpoints.
It is monotone, hence nonnegative everywhere on the interval.
Let
\[
 H_E=(s-r)\max(V(r)^2,V(s)^2).
\]
When \(b\ge0\), take the secant line
\[
 L(u)=a+b\frac{r+s-u}{rs}.
\]
When \(b<0\), take \(c=(r+s)/2\) and
\[
 L(u)=a+2b/c-bu/c^2.
\]
\begin{lemma}[Affine integral certificate]\label{lem:affine}
In both cases \(0\le V\le L\), and with
\(\alpha=L(r)\), \(\beta=L(s)\),
\begin{equation}\label{eq:affine}
 \int_r^sV(u)^2\dd u
 \le\min\left(H_E,\frac{s-r}{3}
                       (\alpha^2+\alpha\beta+\beta^2)\right).
\end{equation}
Every quantity on the right is rational.
\end{lemma}
\begin{proof}
For \(b\ge0\),
\[
 \frac{r+s-u}{rs}-\frac1u=\frac{(u-r)(s-u)}{rsu}\ge0.
\]
For \(b<0\),
\[
 L(u)-V(u)=\frac{(-b)(u-c)^2}{uc^2}\ge0.
\]
Square the inequalities, and integrate the square of the affine
function using \(u=r+t(s-r)\). Its integral is
\((s-r)(\alpha^2+\alpha\beta+\beta^2)/3\).
The endpoint bound \(H_E\) is also valid by monotonicity. The smaller
of two upper bounds for the same integral remains an upper bound.
\end{proof}

If the first interval is \([0,s]\) with verified \(b=0\), its
integral is \(a^2s\); no division by zero is performed.
At fixed integer scale \(S>0\), replace an interval upper bound
\(H\) by \(\lceil SH\rceil/S\). The error is less than \(1/S\)
per interval. At a fixed partition and scale,
\(\lceil S\min(H_E,H_A)\rceil\le\lceil SH_E\rceil\), so this
improvement cannot worsen an existing endpoint certificate.

A strict integral margin guarantees some finite partition and
sufficiently fine rounding scale will succeed. It does not guarantee
that indefinitely subdividing at a fixed scale succeeds: the
accumulated rounding budget depends on the number of intervals.
Nor does a handful of successful leaves certify the remaining leaves.

Apply the envelope and integral bounds to the four direction families
and each of the three other types. Reflections and coordinate exchange
preserve the major-coordinate parameter, so no angular or arc-length
Jacobian is inserted. If the resulting twelve integral upper bounds
sum to \(H^+\), then throughout the root box
\begin{equation}\label{eq:certRisk}
 \mathcal R_r(F;x)\le\frac{D_*}{24}H^+.
\end{equation}
For unscaled fields followed by time contraction by \(s_0\), the
coefficient is \(D_*s_0^2/24\). Reusing a bound on symmetry images
requires verified invariance of the relevant deterministic fields,
or separately transformed fields. Symmetry of the random output is
neither assumed nor required.

\subsection{Capacity: axes, nonaxes and the infinite tail}
For a fixed nonaxis slope \(q\ne0\), the chord of a cell
\((A,B)\times(C,D)\) at intercept \(v\) is
\[
 \pos{\min\left(B,\max\left(\frac{C-v}{q},\frac{D-v}{q}\right)\right)
 -\max\left(A,\min\left(\frac{C-v}{q},\frac{D-v}{q}\right)\right)}.
\]
It is continuous and piecewise affine in \(v\), with switches among
\(C-qA,C-qB,D-qA,D-qB\). Summing finitely many cell contributions
proves that the maximum capacity is attained at an intercept through
a grid vertex. Thus a finite rational list covers every intercept.

For an axis the assertion about endpoint values is false. A horizontal
line in an open row strip has chord equal to that strip's weighted
row sum, whereas a horizontal grid boundary may have chord zero.
Check all \emph{open strip values} and the actual boundary value zero;
use column sums for the vertical axis. Finer zeroed grid lines either
remove finitely many points of a nonaxis chord or set an entire axis
chord to zero, so a coarse nonnegative envelope remains an upper bound.

For the candidate's height cutoff, all directions with \(h\ge8\)
satisfy
\[
 L_{v,r}(F;x)\le\frac{\|F_r+F_{r+\bar v}\|_\infty}{2h}
 \le\frac{\|F_r+F_{r+\bar v}\|_\infty}{16}.
\]
Check the finite cell maxima for all six unordered type pairs, and
check each primitive direction of heights one through seven by the
appropriate axis or nonaxis rule. Nonnegative accumulation makes a
terminal-field capacity bound valid at earlier times as well.

\begin{proposition}[A finite geometric checking contract]\label{prop:contract}
A successful finite checker must establish all of the following:
\begin{enumerate}
\item The exact field values, boundary convention, phase durations,
types, order, and accepted union are the specified inputs.
The actual accumulated field is bounded by the corresponding
reference prefix at every used phase.
\item Every used root lies in a checked closed box; all required
direction families and type pairs have complete angular coverage.
Every branch is justified from its field and box, and every integral
bound uses sound rational arithmetic and outward rounding.
\item The capacities cover the axis strips, the finite nonaxis list,
and the infinite-height tail. The union area counts overlaps once.
\item The returned rational bounds satisfy
\(U_H<K\), \(U_L<\lambda<1-\rho\), and \(U_R<d<1\).
\end{enumerate}
Then the analytic input inequalities hold with uniform positive
margins \(K-U_H\), \(\lambda-U_L\), and \(d-U_R\).
\end{proposition}
\begin{proof}
The capacity rules cover all intercepts and directions.
Proposition~\ref{prop:envelope}, branch induction and
Lemma~\ref{lem:affine} imply \eqref{eq:certRisk} on each closed box.
Nonnegative field domination transfers this to the actual field;
the phase and root-deadline matching ensures it is the correct
field. The finite rational maxima or a common verified threshold
give one uniform bound, whose positive difference from the target
is independent of the grid size and root. The separate union-area
calculation supplies the mass. Proposition~\ref{prop:pair} then
transfers the continuum risk to the discrete pair sum.
\end{proof}

Notice that a stored numerical upper bound or a hash of a data file
does not by itself establish any of the branch, coverage or prefix
premises. Hashes identify the input being discussed; the mathematical
checking contract establishes the inequalities about that input.


\section{The numerical application and its exact budget}\label{sec:numerics}

\subsection{Contraction of a fixed input}
For a fixed \(s_0\in(0,1)\), define a new schedule by multiplying
every duration by \(s_0\), while preserving masks, cells and order.
Then
\[
 F'_{r,t}=s_0F_{r,t/s_0},\qquad \tau'_r=s_0\tau_r,
 \qquad A'=s_0A.
\]
Chord capacities scale by \(s_0\), and squared-chord risk scales by
\(s_0^2\). These identities follow directly by substitution in the
time integrals. They concern a new deterministic input, rather than
an unproved time-change identity for the old random process.

The adopted values from the preparation and followup records are
listed below. Fractions, rather than their decimal displays, are the
values to be checked. Their arithmetic consistency is rechecked by
the script distributed with this draft; that script does not certify
the geometric field from which the stored bounds were obtained.

\begin{center}
\begin{tabular}{@{}ll@{}}
\toprule
Quantity & Recorded value or choice\\
\midrule
Target \(c\) & \(1501/1000\)\\
Time contraction \(s_0\) & \(249/250\)\\
Scaled input mass \(A_*\) & \(5036868114751900569/3355443200000000000\)\\
Total time \(T\) & \(155699823660621/6250000000\)\\
Designated clock cap \(K\) & \(255/128\)\\
Ordinary capacity cap \(\lambda\) & \(4/5\)\\
Ordinary credit floor \(\rho\) & \(1/8\)\\
Weight exponent \(d\) & \(993/1000\)\\
Virtual exponent \(\gamma\) & \(1993/2000\)\\
Geometric and correction exponents \(\xi,\zeta\) & \(1/1000,\ 1/1000\)\\
Finite-sieve coefficient \(D_*\) & \(76/125\)\\
Spatial grid \(G\); designated directions \(g\) & \(32768;\ 16\)\\
\bottomrule
\end{tabular}
\end{center}

The recorded scaled capacity bounds and their margins are
\begin{align*}
 U_H&=\frac{159305594247}{80000000000},
 &K-U_H&=\frac{69405753}{80000000000},\\
 U_L&=\frac{2535641192953}{3200000000000},
 &\lambda-U_L&=\frac{24358807047}{3200000000000},\\
 U_{\rm tail}&=\frac{294190968417}{400000000000},
 &1-\rho-\lambda&=\frac3{40}.
\end{align*}
All margins are positive. The tail is included in the ordinary
capacity contract, not used as a substitute for the finitely many
low-height directions.

\subsection{Changing the directional coefficient without changing its meaning}
The old coefficient, used here only as a rational number, was
\[
 D_{\rm old}=\left(\sum_{k=1}^{256}\frac1{k^2}
                    +\frac1{257}+\frac1{2\cdot257^2}\right)^{-1}.
\]
For the same geometrically justified unscaled envelope integral,
the old scaled bound was
\[
 r_{\rm old}=\frac{3099987999}{3125000000},\qquad
 \kappa=d-r_{\rm old}=\frac{3137001}{3125000000}.
\]
Changing to the finite-sieve coefficient gives the inherited comparison
\begin{align}
 r_{\rm new}&=\frac{D_*}{D_{\rm old}}r_{\rm old}
       <d-\frac{7\kappa}{8},\label{eq:risktransfer}\\
 b_{\rm acc}&=d-\frac\kappa2
       =\frac{6203112999}{6250000000}=0.99249807984.
       \label{eq:bacc}
\end{align}
Thus \(b_{\rm acc}-r_{\rm new}>3\kappa/8>0\) remains for finite
subdivision and rounding. The coefficient multiplying the twelve
unscaled direction/type integrals is exactly
\[
 \frac{D_*s_0^2}{24}=\frac{392673}{15625000}.
\]
This comparison is conditional on the old bound describing the same
valid envelope. It is not a proof that every saved branch partition
is accepted by the new checker. Indeed, increasing \(D\) may make
a previously tight endpoint bound fail; the affine improvement or a
refinement must then be checked for that actual leaf.

There is a further distinction between continuum acceptance and the
discrete pair bound used by the process. If all used roots satisfy
\(\mathcal R_r\le b_{\rm acc}\), choose
\begin{equation}\label{eq:adiscrete}
 a=\frac{b_{\rm acc}+d}{2}=0.99274903992<d.
\end{equation}
Proposition~\ref{prop:pair} gives
\(\mathcal B_n^{\Gamma^c}\le b_{\rm acc}\log n+C\).
For \(\log n\ge C/(a-b_{\rm acc})\), this implies the required
\(\mathcal B_n^{\Gamma^c}\le a\log n\).
One must not silently absorb the additive error while retaining the
same coefficient \(b_{\rm acc}\).

\subsection{Moment choices and the fixed loss budget}
The exponent conditions hold because
\[
 1-d-2\xi-\zeta=\frac1{250}>0,\qquad
 \gamma-d=\frac7{2000}>0.
\]
For example \(m=4096\) makes the exponents in
\eqref{eq:etatail} and \eqref{eq:virtualtail}, before point-law
transport, equal to \(-13.377\) and \(-12.328\).
These are illustrative fixed choices, not a uniform statement about
the hidden constants for every order.

For the tag bound, \(C_{\rm line}=4G=131072\). The recorded fixed
budget \(\varepsilon=10^{-5}\), with one third assigned to tags,
can use
\[
 e=\frac{1-\rho-\lambda}{4TC_{\rm line}},\qquad
 b=348294609696,\quad k=116,\quad
 m_{\rm tag}=bk+3=40402174724739.
\]
The rational witnesses
\(b(e/2)\ge1\) and
\(2^{-k}\le (1/300000)e^2/(64C_{\rm line})\) imply the desired
small main term in \eqref{eq:tagprob}, by Bernoulli's inequality
\((1+e/2)^b\ge2\). Its enormous but fixed degree is permitted by
the proof: a general fixed-order moment argument is used before
substituting this integer. No such power is expanded into individual
terms in the arithmetic check. This choice supplies neither a useful
explicit \(n_0\) nor an estimate for practical simulation size.

\begin{corollary}[Candidate application, conditional on full input checking]\label{cor:target}
Suppose the candidate accepted schedule has mass \(A_*\), the stated
capacity bounds, and a valid geometric certificate giving risk at
most \(b_{\rm acc}\) at every used root's deadline. Suppose the
accepted unions, durations and reference prefixes have the matching
specified in Proposition~\ref{prop:contract}. Then the analytic
reduction yields \eqref{eq:target}.
\end{corollary}
\begin{proof}
The finite sieve, positive margins and \eqref{eq:adiscrete} supply
the geometric hypotheses; the displayed parameters satisfy
\eqref{eq:exponents}. Theorem~\ref{thm:bridge} applies.
Alternatively, retain just one error budget: choose \(n\) large
enough that \(A_n\ge A_*-\Delta/4\) and the total early-retirement
point-law probability is at most \(10^{-5}\). Proposition~\ref{prop:loss}
and the strict inequality in Section~\ref{sec:overview} give
\(\E|S_n|>cn\). A legal realization then satisfies
\(1501n\le1000|S_n|\).
\end{proof}

A separately verified pruning loss \(P\le\Delta/4\) can also be
paid in the same ledger, provided the resulting schedule continues
to satisfy the required geometric and matching conditions. With
grid error \(\Delta/4\),
\[
 (1-10^{-5})(A_* -\Delta/4-\Delta/4)-c
 =0.000036826382327236234293878078460693359375>0.
\]
The two occurrences of \(\Delta/4\) pay different costs; neither
is reused. This is a budget lemma, not a claim that a pruned input
has actually been constructed or checked for this manuscript.

\section{Verification status, correspondence and provenance}\label{sec:status}

\subsection{What this document adds}
This document is an English mathematical exposition of the preparation
notes and their two followups~\cite{preparation,followup,second},
organized around the exact loss identity.
It adopts finite inclusion-exclusion in place of the older infinite
Mobius argument, affine rational integral bounds in place of mandatory
logarithmic integration, and the linewise cell count in the tag proof.
The arguments are written out for a reader, rather than replaced by
Lean declarations. They remain a research proof draft under review.
Editing and arithmetic checks are not independent verification of the
entire stochastic argument or the concrete numerical certificate.

The paper is self-contained in the sense that its mathematical
definitions, proof steps and the checking contract are stated here.
It is not a standalone accepted numerical certificate: the large
reference field and all geometric leaves are not embedded in the
text or in the small companion data file. The numerical corollary
retains those explicit premises.

\subsection{The fixed Lean snapshot}
The comparison here is with the public snapshot~\cite{formal} whose clean-history
commit is \texttt{282b636}, based on the audit of 10 November 2026,
00:22:46 JST. It contains 3,242 original Lean modules and reports
839,734 source lines. This count measures source inventory, not the
percentage of the final theorem proved. The fixed snapshot records:
\begin{itemize}
\item actual one-phase exponential estimates and uniform power tails;
\item actual frozen-schedule clock-cause smallness;
\item an output theorem combining the five causes under explicit
concrete-input hypotheses;
\item 99 of 112 low-height capacity rows, and the specified nonaxis
capacity table with 28 of 28 rows;
\item 20 of 80 canonical leaves at the single benchmark phase 15301,
with 80 of 320 symmetry images.
\end{itemize}
The last count is not 25 percent of the complete geometric input.
The full input has 15,348 phases and 155,402 leaf records.

The principal formal endpoint is
\begin{center}
\small\path{FrozenCapacity.exists_uniform_accepted_frozen_output_of_inputs}
\end{center}
in \path{NoThreeLine/FrozenFiveCauseOutput.lean}.
Its hypotheses include all 112 accepted rows; a fixed number of
checked accepted phases; their matching to the frozen reference with
strictly increasing indices; their duration bounds and fixed total
time; actual grid input mass above \(1501/1000\); and the actual
normalized pair-sum bound at every used root's deadline, with fixed
\(0\le a<d<\gamma<1\) and \(0<\zeta<1-d\).
It does not assert that these inputs have been constructed for all
sufficiently large \(n\).

Two intermediate entry points are
\path{NoThreeLine/ProtectedCellOnePhasePowerTail.lean} and
\path{NoThreeLine/FrozenClockCauseSmall.lean}.
The general expectation-to-output interface is in
\path{NoThreeLine/ScheduleFiveCauseOutput.lean}.
The paper's explanatory grouping is not a claim that every displayed
lemma is a verbatim translation of one Lean declaration. The formal
source, with its actual hypotheses, determines the formal scope.

\subsection{The remaining numerical and mathematical obligations}
To turn Corollary~\ref{cor:target} into an unconditional claim, one
must complete a single compatible chain of evidence:
\begin{enumerate}
\item Certify every adopted root region with its correct field,
parity labels, angular coverage and symmetry images. If using the
new finite-sieve coefficient, establish acceptance under that
coefficient, not just under the older one.
\item Establish chronological phase matching and all cumulative
reference-prefix bounds. A local accepted rectangle alone is not
a valid complete schedule.
\item Check union mass, capacity maxima, axis strips and direction
tails for that same input. Complete the remaining finite row
obligations if using the frozen Lean interface.
\item Transfer those continuum bounds to the actual grid input and
actual root-deadline pair sum with one sufficiently large threshold
uniform in all used roots. The additive-error step in
\eqref{eq:adiscrete} is part of this obligation.
\end{enumerate}
The probabilistic proof also deserves direct mathematical examination,
especially common stopping of mixed hierarchies, whole-event energy,
conditional restart at the first line birth, and the exceptional
configurations in the short-star expansion. The recorded AI reviews
and partial formalization are evidence at their stated scopes; they
do not convert the remaining premises into facts.

\subsection{Companion files and what was checked for this edition}
The editable source is \path{main.tex} with its section files.
The small file \path{data/parameters.json} records the exact numerical
choices and the finite-sieve factors. Running
\path{scripts/check_arithmetic.py} checks the mass and capacity
margins, the coefficient transfer, the discrete-risk slack,
the two fixed loss budgets, the finite product and coprimality
conditions, and the tag-order witnesses. Its report is a verification
of those finite arithmetic statements, not an execution of the full
geometric checker or a check of the whole proof.

\path{SOURCE_MAP.json} identifies the source notes and the frozen
public modules used in preparing this edition. Source identifiers
and hashes serve reproducibility; their role is distinct from the
soundness arguments in Section~\ref{sec:certificates}.
The PDF has been typeset and visually checked. This edition performs
no new large angular computation and no full Lean rebuild.

\subsection{AI use and licensing}
This work has involved substantial use of OpenAI ChatGPT/Codex and
AI reviews supplied by the project maintainer. AI systems contributed
mathematical proposals and derivations, numerical programs, exact
arithmetic checks, Lean definitions and proofs, build orchestration,
review, translation and exposition. Separate AI contexts reviewed
parts of the earlier work. This English edition was assembled and
checked by the assistant in this task; no new independent human
review or full independent proof audit is claimed.

The human maintainer chose the direction and requested experiments
and revisions as a hobby project. This disclosure does not attribute
every mathematical or programming step to human authorship. No named
author, affiliation or personal contact address is supplied. The
repository's hosting account is public. The original project material
and this draft are distributed under the MIT license supplied in
\path{LICENSE}; retain its notices when redistributing. No claim of
anonymous hosting or of formal completion is made.

\begin{thebibliography}{9}
\bibitem{preparation}
Project preparation notes, \emph{Integrated proof draft preceding
Lean formalization}, 8--9 November 2026 (Japanese).
\href{https://github.com/agolkjaoi/no-three-in-line-lean/blob/282b636c160d049e3ef4d82f619fd44d98f8cf1d/docs/research-notes/preparation/PROOF_FOR_LEAN_JA.txt}{Public fixed-version text}.
This draft uses its process, probability bridge, uniform geometry,
and stopping/risk arguments.
\bibitem{followup}
Project first followup, 9 November 2026 (Japanese).
\href{https://github.com/agolkjaoi/no-three-in-line-lean/blob/282b636c160d049e3ef4d82f619fd44d98f8cf1d/docs/research-notes/first-followup/FOLLOWUP_FOR_REVIEW_JA.txt}{Public fixed-version text}.
Source of the remaining-exposure identity, fixed loss budget,
finite height sum and strengthened capacity contract.
\bibitem{second}
Project second followup, 9 November 2026 (Japanese).
\href{https://github.com/agolkjaoi/no-three-in-line-lean/blob/282b636c160d049e3ef4d82f619fd44d98f8cf1d/docs/research-notes/second-followup/SECOND_FOLLOWUP_FOR_REVIEW_JA.txt}{Public fixed-version text}.
Source of the finite sieve, affine integral bounds and linewise
cell count.
\bibitem{formal}
\emph{no-three-in-line-lean}, partial formalization snapshot,
10 November 2026.
\href{https://github.com/agolkjaoi/no-three-in-line-lean/tree/282b636c160d049e3ef4d82f619fd44d98f8cf1d}{Source snapshot};
\href{https://github.com/agolkjaoi/no-three-in-line-lean/blob/282b636c160d049e3ef4d82f619fd44d98f8cf1d/STATUS.md}{status and remaining obligations}.
\end{thebibliography}


\end{document}
